\documentclass[10pt,a4paper]{article}
\usepackage[margin=2.2cm]{geometry}
\usepackage{fontspec}
\setmonofont{DejaVu Sans Mono}[Scale=0.84]
\newfontfamily\unifont{DejaVu Serif}[Scale=0.9]
\usepackage{amsmath,amssymb,mathtools}
\setcounter{MaxMatrixCols}{20}
\usepackage{longtable,booktabs,array}
\usepackage{fancyvrb}
\usepackage[dvipsnames]{xcolor}
\usepackage[colorlinks=true,linkcolor=NavyBlue,urlcolor=NavyBlue]{hyperref}
\usepackage{microtype}
\usepackage{needspace}
\usepackage{tocloft}\setlength{\cftbeforesecskip}{0.45em}
\setlength{\parindent}{0pt}\setlength{\parskip}{0.55em}
\setlength{\emergencystretch}{8em}
\makeatletter
\renewcommand\section{\@startsection{section}{1}{\z@}{-3.5ex \@plus -1ex \@minus -.2ex}{2.3ex \@plus.2ex}{\normalfont\Large\bfseries\raggedright}}
\renewcommand\subsection{\@startsection{subsection}{2}{\z@}{-3.25ex\@plus -1ex \@minus -.2ex}{1.5ex \@plus .2ex}{\normalfont\large\bfseries\raggedright}}
\renewcommand\subsubsection{\@startsection{subsubsection}{3}{\z@}{-3.25ex\@plus -1ex \@minus -.2ex}{1.5ex \@plus .2ex}{\normalfont\normalsize\bfseries\raggedright}}
\makeatother
\renewcommand{\arraystretch}{1.15}
\allowdisplaybreaks
\begin{document}
\title{dirac16complex: provenance\\[0.4em]\large Sources, decisions, derivations and the verification record of the Dirac equation and the Kohn–Sham theories of the 16-component field dirac16complex — sections 7–10 of the notebook 7+1.nb}
\author{Repository: github.com/once-ere/7plus1 — prepared for Nash\\ Notebook: 7+1.nb (Wolfram Language 15.0.1)\\ Document generated by tools/make\_docs.py on 2026-10-28}
\date{}
\maketitle
\begin{abstract}
This record states what was asked, what was used (the first part of 7+1.nb and the literature), every decision taken with the equation behind it, how each result was derived, and how it was checked: the 74 checks of sections 7–9 of 7+1.nb with their results, the 90 boolean checks of the independent Python program, the comparison of every printed number, the program versions and the SHA-256 hashes of the files.
\end{abstract}
\tableofcontents
\Needspace{12\baselineskip}\section*{1. The task}\addcontentsline{toc}{section}{1. The task}
This stage of the project asked for the following (summarized; the original wording is the author’s private working material and is not reproduced):

\begin{enumerate}
\item Back up 7+1.nb, then append to it (the first part of the notebook stays as it is).
\item A valid and checked Dirac equation on the real 8-dimensional spacetime with one time and seven space dimensions, for a 16-component quantized fermion field named dirac16complex whose components are complex-valued Grassmann (anticommuting) operators.
\item Eight 16×16 Dirac matrices taken from an appropriate set of anticommuting real matrices of 7+1.nb; they MUST be purely real.
\item Chiral projections that write the 16-component field equation as two 8-component field equations, coupled when the mass is not zero.
\item A 16-component Kohn–Sham theory of dirac16complex coupled to itself.
\item A chiral version of that Kohn–Sham theory, using the chiral projections.
\item Complete, precise and checked documentation; for every documentation file a companion provenance file and a symbol file, each in Markdown and LaTeX, the LaTeX compiled to PDF; no document may refer the reader to another document, and every equation used must be written in the document that uses it.
\item Push every completed stage to the repository github.com/once-ere/7plus1 (the repository was renamed from 7+1 to 7plus1 during this stage), respecting the repository’s .gitignore.
\item Use only the Wolfram front end WolframNB.exe of Wolfram 15.0.1, installed in \texttt{C:\textbackslash{}\allowbreak{}Program Files\textbackslash{}\allowbreak{}Wolfram Research\textbackslash{}\allowbreak{}Wolfram\textbackslash{}\allowbreak{}15.0.1}.
\end{enumerate}
The coordinates are those of 7+1.nb and were not renamed: $X=(x_0,x_1,\dots,x_7)$ with $x_0$ the ordinary time (the direction of the octonion unit $e_0=1$), $x_1,x_2,x_3$ ordinary 3-space and $x_4,\dots,x_7$ four auxiliary space-like dimensions. How each requirement was met:

{\footnotesize
\begin{longtable}{>{\raggedright\arraybackslash}p{\dimexpr 0.0500\linewidth-2\tabcolsep\relax}>{\raggedright\arraybackslash}p{\dimexpr 0.7500\linewidth-2\tabcolsep\relax}>{\raggedright\arraybackslash}p{\dimexpr 0.2000\linewidth-2\tabcolsep\relax}}
\toprule
\textbf{\#} & \textbf{how it was met} & \textbf{where in 7+1.nb} \\
\midrule\endfirsthead
\toprule
\textbf{\#} & \textbf{how it was met} & \textbf{where in 7+1.nb} \\
\midrule\endhead
\endfoot
\bottomrule\endlastfoot
1 & the evaluated notebook of the first stage was copied to backups/7+1\_2026-10-26\_before\_dirac16complex.nb (section 2.1); sections 0–6 of the new notebook are cell for cell identical to it (section 4.3), except for the cells of the signature explanation of 27 October 2026 (section 3.16) & sections 0–6 unchanged (but for 3.16) \\
2 & $(i\Gamma^A\partial_A-m)\psi=0$ for \texttt{dirac16complex = Table[ψ16[a] @@ X, \{a, 1, 16\}]}; Klein–Gordon, Hamiltonian, plane waves, Grassmann algebra, mode expansion, canonical anticommutators, charge conjugation and Lorentz covariance (section 3.2–3.6) & 7.1–7.9 \\
3 & $\Gamma_0=\left(\begin{smallmatrix}0&1_8\\1_8&0\end{smallmatrix}\right)$, $\Gamma_i=\left(\begin{smallmatrix}\tau^{O}_i&0\\0&-\tau^{O}_i\end{smallmatrix}\right)$: the real generators \texttt{Γ17[A]} of the first stage; every entry $0$ or $\pm1$ (section 3.1) & 7.1 \\
4 & complex chirality $\chi=i\Gamma_0\cdots\Gamma_7$ (no real one exists), $P_{L,R}=\frac12(1\mp\chi)$, and the two coupled equations $(i\partial_0\pm\tau^{O}_i\partial_i)\phi_{L,R}=m\,\phi_{R,L}$ (section 3.7) & 7.10–7.11 \\
5 & self-coupling through scalar, pseudoscalar, vector and axial-vector contact interactions; Hohenberg–Kohn theorem; Kohn–Sham equations; exchange (Wick), Fierz; uniform matter (sections 3.8–3.11, 3.13) & 8.1–8.7 \\
6 & chiral densities $\kappa$, $j_L$, $j_R$; complex Kohn–Sham mass; 8+8-component Kohn–Sham equations; chiral symmetry; the chiral vacuum with the Dirac sea (sections 3.12, 3.14) & 9.1–9.6 \\
7 & the documentation of this stage (a main document, a symbol dictionary and this record, each .md/.tex/.pdf), generated by tools/make\_docs.py (section 5) & — \\
8 & commits pushed to github.com/once-ere/7plus1 by the Wolfram job \texttt{tools/\allowbreak{}git\_\allowbreak{}push.wl} (section 5) & — \\
9 & every evaluation of 7+1.nb was made by that front end (section 4.1) & — \\*
\end{longtable}}
\Needspace{12\baselineskip}\section*{2. Inputs}\addcontentsline{toc}{section}{2. Inputs}
\Needspace{8\baselineskip}\subsection*{2.1 The notebook 7+1.nb of the first stage}\addcontentsline{toc}{subsection}{2.1 The notebook 7+1.nb of the first stage}
Before anything was appended, the evaluated notebook of the first stage (128 checks, all True) was copied, in the workspace \texttt{D:\textbackslash{}\allowbreak{}Developer\textbackslash{}\allowbreak{}github\textbackslash{}\allowbreak{}7+1}, to the file below. The folder backups/ is excluded from the repository by its .gitignore, so the copy stays in the workspace only.

\begin{Verbatim}[fontsize=\small,frame=single,framesep=2mm]
backups/7+1_2026-10-26_before_dirac16complex.nb
306963 bytes
SHA-256 f7abe7279b5540e914de876dc055c199bd491b9c2d9522bf28dbc23689e237ce
\end{Verbatim}
Sections 7–10 use the following objects of sections 1–5 of 7+1.nb, unchanged (Mathematica names in typewriter). Coordinates, metric and derivative operators:
\begin{align*}
&X=(x_0,x_1,\dots,x_7)\ (\texttt{X}),\qquad \eta=\mathrm{diag}(+1,-1,-1,-1,-1,-1,-1,-1)\ (\texttt{η17}),\qquad \partial_A=\partial/\partial x_A, \\
&\Gamma^A\partial_A\psi\ (\texttt{DiracOp17}),\qquad \Box\psi=\eta^{AB}\partial_A\partial_B\psi=\partial_0^2\psi-\partial_1^2\psi-\dots-\partial_7^2\psi\ (\texttt{waveOp17}).
\end{align*}
The real 8×8 octonion matrices \texttt{τO[i]} ($\tau^{O}_0=1_8$), obtained in the first stage from the split-octonion matrices $\tau_A$ of 4+4.nb as $\tau^{O}_i=\tau_i$ ($i=1,2,3$) and $\tau^{O}_k=\tau_k\,\sigma$ ($k=4,\dots,7$), $\sigma=\tau_1\tau_2\tau_3=\left(\begin{smallmatrix}0&1_4\\1_4&0\end{smallmatrix}\right)$:
\begin{equation*}
\tau^{O}_1=\begin{pmatrix}0 & 0 & 0 & 0 & 0 & 0 & 0 & 1 \\ 0 & 0 & 0 & 0 & 0 & 0 & 1 & 0 \\ 0 & 0 & 0 & 0 & 0 & -1 & 0 & 0 \\ 0 & 0 & 0 & 0 & -1 & 0 & 0 & 0 \\ 0 & 0 & 0 & 1 & 0 & 0 & 0 & 0 \\ 0 & 0 & 1 & 0 & 0 & 0 & 0 & 0 \\ 0 & -1 & 0 & 0 & 0 & 0 & 0 & 0 \\ -1 & 0 & 0 & 0 & 0 & 0 & 0 & 0\end{pmatrix},\quad \tau^{O}_2=\begin{pmatrix}0 & 0 & 0 & 0 & 0 & 0 & -1 & 0 \\ 0 & 0 & 0 & 0 & 0 & 0 & 0 & 1 \\ 0 & 0 & 0 & 0 & 1 & 0 & 0 & 0 \\ 0 & 0 & 0 & 0 & 0 & -1 & 0 & 0 \\ 0 & 0 & -1 & 0 & 0 & 0 & 0 & 0 \\ 0 & 0 & 0 & 1 & 0 & 0 & 0 & 0 \\ 1 & 0 & 0 & 0 & 0 & 0 & 0 & 0 \\ 0 & -1 & 0 & 0 & 0 & 0 & 0 & 0\end{pmatrix},
\end{equation*}
\begin{equation*}
\tau^{O}_3=\begin{pmatrix}0 & 0 & 0 & 0 & 0 & 1 & 0 & 0 \\ 0 & 0 & 0 & 0 & -1 & 0 & 0 & 0 \\ 0 & 0 & 0 & 0 & 0 & 0 & 0 & 1 \\ 0 & 0 & 0 & 0 & 0 & 0 & -1 & 0 \\ 0 & 1 & 0 & 0 & 0 & 0 & 0 & 0 \\ -1 & 0 & 0 & 0 & 0 & 0 & 0 & 0 \\ 0 & 0 & 0 & 1 & 0 & 0 & 0 & 0 \\ 0 & 0 & -1 & 0 & 0 & 0 & 0 & 0\end{pmatrix},\quad \tau^{O}_4=\begin{pmatrix}0 & 1 & 0 & 0 & 0 & 0 & 0 & 0 \\ -1 & 0 & 0 & 0 & 0 & 0 & 0 & 0 \\ 0 & 0 & 0 & -1 & 0 & 0 & 0 & 0 \\ 0 & 0 & 1 & 0 & 0 & 0 & 0 & 0 \\ 0 & 0 & 0 & 0 & 0 & -1 & 0 & 0 \\ 0 & 0 & 0 & 0 & 1 & 0 & 0 & 0 \\ 0 & 0 & 0 & 0 & 0 & 0 & 0 & 1 \\ 0 & 0 & 0 & 0 & 0 & 0 & -1 & 0\end{pmatrix},
\end{equation*}
\begin{equation*}
\tau^{O}_5=\begin{pmatrix}0 & 0 & -1 & 0 & 0 & 0 & 0 & 0 \\ 0 & 0 & 0 & -1 & 0 & 0 & 0 & 0 \\ 1 & 0 & 0 & 0 & 0 & 0 & 0 & 0 \\ 0 & 1 & 0 & 0 & 0 & 0 & 0 & 0 \\ 0 & 0 & 0 & 0 & 0 & 0 & 1 & 0 \\ 0 & 0 & 0 & 0 & 0 & 0 & 0 & 1 \\ 0 & 0 & 0 & 0 & -1 & 0 & 0 & 0 \\ 0 & 0 & 0 & 0 & 0 & -1 & 0 & 0\end{pmatrix},\quad \tau^{O}_6=\begin{pmatrix}0 & 0 & 0 & -1 & 0 & 0 & 0 & 0 \\ 0 & 0 & 1 & 0 & 0 & 0 & 0 & 0 \\ 0 & -1 & 0 & 0 & 0 & 0 & 0 & 0 \\ 1 & 0 & 0 & 0 & 0 & 0 & 0 & 0 \\ 0 & 0 & 0 & 0 & 0 & 0 & 0 & 1 \\ 0 & 0 & 0 & 0 & 0 & 0 & -1 & 0 \\ 0 & 0 & 0 & 0 & 0 & 1 & 0 & 0 \\ 0 & 0 & 0 & 0 & -1 & 0 & 0 & 0\end{pmatrix},
\end{equation*}
\begin{equation*}
\tau^{O}_7=\begin{pmatrix}0 & 0 & 0 & 0 & -1 & 0 & 0 & 0 \\ 0 & 0 & 0 & 0 & 0 & -1 & 0 & 0 \\ 0 & 0 & 0 & 0 & 0 & 0 & -1 & 0 \\ 0 & 0 & 0 & 0 & 0 & 0 & 0 & -1 \\ 1 & 0 & 0 & 0 & 0 & 0 & 0 & 0 \\ 0 & 1 & 0 & 0 & 0 & 0 & 0 & 0 \\ 0 & 0 & 1 & 0 & 0 & 0 & 0 & 0 \\ 0 & 0 & 0 & 1 & 0 & 0 & 0 & 0\end{pmatrix};
\end{equation*}
the split matrices used for them are
\begin{equation*}
\tau_4=\begin{pmatrix}0 & 0 & 0 & 0 & 0 & 1 & 0 & 0 \\ 0 & 0 & 0 & 0 & -1 & 0 & 0 & 0 \\ 0 & 0 & 0 & 0 & 0 & 0 & 0 & -1 \\ 0 & 0 & 0 & 0 & 0 & 0 & 1 & 0 \\ 0 & -1 & 0 & 0 & 0 & 0 & 0 & 0 \\ 1 & 0 & 0 & 0 & 0 & 0 & 0 & 0 \\ 0 & 0 & 0 & 1 & 0 & 0 & 0 & 0 \\ 0 & 0 & -1 & 0 & 0 & 0 & 0 & 0\end{pmatrix},\quad \tau_5=\begin{pmatrix}0 & 0 & 0 & 0 & 0 & 0 & -1 & 0 \\ 0 & 0 & 0 & 0 & 0 & 0 & 0 & -1 \\ 0 & 0 & 0 & 0 & 1 & 0 & 0 & 0 \\ 0 & 0 & 0 & 0 & 0 & 1 & 0 & 0 \\ 0 & 0 & 1 & 0 & 0 & 0 & 0 & 0 \\ 0 & 0 & 0 & 1 & 0 & 0 & 0 & 0 \\ -1 & 0 & 0 & 0 & 0 & 0 & 0 & 0 \\ 0 & -1 & 0 & 0 & 0 & 0 & 0 & 0\end{pmatrix},
\end{equation*}
\begin{equation*}
\tau_6=\begin{pmatrix}0 & 0 & 0 & 0 & 0 & 0 & 0 & -1 \\ 0 & 0 & 0 & 0 & 0 & 0 & 1 & 0 \\ 0 & 0 & 0 & 0 & 0 & -1 & 0 & 0 \\ 0 & 0 & 0 & 0 & 1 & 0 & 0 & 0 \\ 0 & 0 & 0 & 1 & 0 & 0 & 0 & 0 \\ 0 & 0 & -1 & 0 & 0 & 0 & 0 & 0 \\ 0 & 1 & 0 & 0 & 0 & 0 & 0 & 0 \\ -1 & 0 & 0 & 0 & 0 & 0 & 0 & 0\end{pmatrix},\quad \tau_7=\begin{pmatrix}-1 & 0 & 0 & 0 & 0 & 0 & 0 & 0 \\ 0 & -1 & 0 & 0 & 0 & 0 & 0 & 0 \\ 0 & 0 & -1 & 0 & 0 & 0 & 0 & 0 \\ 0 & 0 & 0 & -1 & 0 & 0 & 0 & 0 \\ 0 & 0 & 0 & 0 & 1 & 0 & 0 & 0 \\ 0 & 0 & 0 & 0 & 0 & 1 & 0 & 0 \\ 0 & 0 & 0 & 0 & 0 & 0 & 1 & 0 \\ 0 & 0 & 0 & 0 & 0 & 0 & 0 & 1\end{pmatrix},
\end{equation*}
and $\tau_i=\tau^{O}_i$ for $i=1,2,3$. They satisfy
\begin{equation*}
(\tau^{O}_i)^{\mathsf T}=-\tau^{O}_i,\qquad \tau^{O}_i\tau^{O}_j+\tau^{O}_j\tau^{O}_i=-2\,\delta_{ij}\,1_8,\qquad \tau^{O}_1\cdots\tau^{O}_7=1_8 .
\end{equation*}
The real 16×16 generators \texttt{Γ17[A]} of the Clifford algebra of 7+1 spacetime, their products and the objects built from them:
\begin{align*}
\Gamma_0&=\begin{pmatrix}0&1_8\\1_8&0\end{pmatrix},\qquad \Gamma_i=\begin{pmatrix}\tau^{O}_i&0\\0&-\tau^{O}_i\end{pmatrix},\qquad \Gamma_A\Gamma_B+\Gamma_B\Gamma_A=2\,\eta_{AB}\,1_{16},\qquad \Gamma^A=\eta^{AA}\Gamma_A\ (\texttt{Γ17up[A]}), \\
\Gamma_8&\equiv\omega=\Gamma_0\Gamma_1\cdots\Gamma_7\ (\texttt{Γ17[8]}),\qquad \chi=i\,\omega\ (\texttt{χ17}),\qquad P_{L,R}=\tfrac12(1_{16}\mp\chi)\ (\texttt{PL17},\ \texttt{PR17}), \\
\Sigma_{AB}&=\tfrac14(\Gamma_A\Gamma_B-\Gamma_B\Gamma_A)\ (\texttt{Σ17[A, B]}),\qquad \Gamma_K=\Gamma_{A_1}\cdots\Gamma_{A_k}\ (K=\lbrace A_1<\dots<A_k\rbrace;\ \texttt{products256}), \\
\mathcal B(\beta_{\rm b})&=\cosh\tfrac{\beta_{\rm b}}2\,1_{16}+\sinh\tfrac{\beta_{\rm b}}2\,\Gamma_0\Gamma_1\ (\texttt{Sβ}),\qquad L_{BK}=\tfrac1{16}\mathrm{tr}\big(\eta_{BB}\Gamma_B\,\mathcal B\,\Gamma_K\,\mathcal B^{-1}\big)\ (\texttt{Λboost}).
\end{align*}
In terms of the 16×16 Dirac matrices of 4+4.nb, $T_A=\left(\begin{smallmatrix}0&\bar\tau_A\\\tau_A&0\end{smallmatrix}\right)$ with $\bar\tau_0=1_8$, $\bar\tau_A=\sigma\tau_A^{\mathsf T}\sigma$ and $T_AT_B+T_BT_A=2\eta^{(4,4)}_{AB}1_{16}$, $\eta^{(4,4)}=\mathrm{diag}(+1,+1,+1,+1,-1,-1,-1,-1)$, the first stage obtained $\Gamma_0=T_0$, $\Gamma_i=T_0T_i$ ($i=1,2,3$) and $\Gamma_k=T_0T_kT_4T_5T_6T_7$ ($k=4,\dots,7$).

\Needspace{8\baselineskip}\subsection*{2.2 Literature, and the equation taken from each source}\addcontentsline{toc}{subsection}{2.2 Literature, and the equation taken from each source}
\begin{itemize}
\item P. A. M. Dirac, Proc. R. Soc. Lond. A 117 (1928) 610 — the first-order wave equation $(i\gamma^\mu\partial_\mu-m)\psi=0$ with $\gamma^\mu\gamma^\nu+\gamma^\nu\gamma^\mu=2\eta^{\mu\nu}$; here written in 7+1 dimensions with the eight $\Gamma^A$.
\item P. Hohenberg and W. Kohn, Phys. Rev. 136 (1964) B864 — for a non-degenerate ground state the density determines the external potential up to a constant, and $E[v]=\min_n\big(F[n]+\int v\,n\big)$; the proof by the two variational inequalities is used in section 3.10.
\item W. Kohn and L. J. Sham, Phys. Rev. 140 (1965) A1133 — a non-interacting reference system with the same density: $\big[-\tfrac12\nabla^2+v+v_H+v_{xc}\big]\varphi_k=\varepsilon_k\varphi_k$, $n=\sum_k|\varphi_k|^2$, $v_{xc}=\delta E_{xc}/\delta n$ (atomic units); here its relativistic form.
\item A. K. Rajagopal and J. Callaway, Phys. Rev. B 7 (1973) 1912; A. H. MacDonald and S. H. Vosko, J. Phys. C 12 (1979) 2977 — relativistic density-functional theory: the energy is a functional of the current $j^A$, and the Kohn–Sham equations have the Dirac form $\big[\alpha_i(-i\partial_i+\mathcal V_i)+\beta m+\mathcal V_0\big]\varphi_k=\varepsilon_k\varphi_k$ with the effective potential $\mathcal V_A=\delta(E-T_s)/\delta j^A$.
\item M. Levy, Proc. Natl. Acad. Sci. USA 76 (1979) 6062; E. H. Lieb, Int. J. Quantum Chem. 24 (1983) 243 — the universal functional by constrained search, $F[n]=\min_{\Psi\to n}\langle\Psi|\hat T+\hat W|\Psi\rangle$, and Lieb’s form $F[n]=\sup_v\big(E[v]-\int v\,n\big)$.
\item R. M. Dreizler and E. K. U. Gross, \emph{Density Functional Theory}, Springer (1990), ch. 8; E. Engel and R. M. Dreizler, \emph{Density Functional Theory: An Advanced Course}, Springer (2011) — the no-sea approximation: the densities are sums over the occupied positive-energy Kohn–Sham states only, $j^A=\sum_{0<\varepsilon_k\le\varepsilon_F}\bar\varphi_k\Gamma^A\varphi_k$.
\item J. D. Walecka, Ann. Phys. 83 (1974) 491; B. D. Serot and J. D. Walecka, Adv. Nucl. Phys. 16 (1986) 1 — relativistic mean field of uniform matter: $m^*=m-\frac{g_\sigma^2}{m_\sigma^2}\rho_S$, $V_0=\frac{g_\omega^2}{m_\omega^2}\rho$, $\varepsilon=\frac{g_\omega^2}{2m_\omega^2}\rho^2+\frac{m_\sigma^2}{2g_\sigma^2}(m-m^*)^2+\varepsilon_{\rm kin}(m^*)$; the contact couplings here correspond to $g_S=-g_\sigma^2/m_\sigma^2$ and $g_V=g_\omega^2/m_\omega^2$.
\item B. A. Nikolaus, T. Hoch and D. G. Madland, Phys. Rev. C 46 (1992) 1757 — relativistic point-coupling (contact) models, with four-fermion terms of the form $-\frac{\alpha_S}{2}(\bar\psi\psi)^2-\frac{\alpha_V}{2}(\bar\psi\gamma_\mu\psi)(\bar\psi\gamma^\mu\psi)$, the 3+1-dimensional analogue of the $g_S$ and $g_V$ terms here.
\item N. M. Hugenholtz and L. Van Hove, Physica 24 (1958) 363 — at saturation (zero pressure $P=\rho\mu-\varepsilon$) the Fermi energy equals the energy per particle, $\varepsilon_F=\mu=E/A$.
\item G. C. Wick, Phys. Rev. 80 (1950) 268 — expectation values in a Slater determinant factorize into contractions: $\langle c_a^\dagger c_c^\dagger c_dc_b\rangle=\varrho_{ba}\varrho_{dc}-\varrho_{da}\varrho_{bc}$, $\varrho_{xy}=\langle c_y^\dagger c_x\rangle$.
\item M. Fierz, Z. Phys. 104 (1937) 553 — rearrangement of four-fermion products from the completeness of the Clifford basis, here $\sum_K(\Gamma_K)_{ab}(\Gamma_K^{-1})_{cd}=16\,\delta_{ad}\delta_{cb}$.
\item Y. Nambu and G. Jona-Lasinio, Phys. Rev. 122 (1961) 345; S. P. Klevansky, Rev. Mod. Phys. 64 (1992) 649 — chirally symmetric contact coupling $G[(\bar\psi\psi)^2+(\bar\psi i\gamma_5\psi)^2]$ with a momentum cutoff; gap equation $M=m-2G\langle\bar\psi\psi\rangle$; for $m=0$ a solution $M\neq0$ exists above a critical coupling (here $G=-g/2$).
\item F. A. Berezin, \emph{The Method of Second Quantization}, Academic Press (1966) — Grassmann algebra, $\theta_k\theta_l=-\theta_l\theta_k$, and the left derivative.
\item H. B. Lawson and M.-L. Michelsohn, \emph{Spin Geometry}, Princeton (1989) — the classification of real Clifford algebras. In the convention used here ($Cl(p,q)$ has $p$ generators squaring to $+1$ and $q$ to $-1$) it gives $Cl(1,7)\cong\mathbb R(16)$ and $Cl(7,1)\cong\mathbb H(8)$; Lawson and Michelsohn define $v^2=-q(v)$, so their $Cl_{7,1}$ is the $Cl(1,7)$ of this record.
\end{itemize}
\Needspace{12\baselineskip}\section*{3. Decisions and derivations}\addcontentsline{toc}{section}{3. Decisions and derivations}
Each subsection states a decision or a derived result, the equations behind it, and the checks of 7+1.nb that confirm it (quoted by their names in section 4.2).

\Needspace{8\baselineskip}\subsection*{3.1 The real Dirac matrices and the sign of the metric}\addcontentsline{toc}{subsection}{3.1 The real Dirac matrices and the sign of the metric}
\textbf{Decision.} The Dirac matrices are the real generators $\Gamma_A$ of section 2.1, with
\begin{equation*}
\Gamma_A\Gamma_B+\Gamma_B\Gamma_A=2\,\eta_{AB}\,1_{16},\qquad \eta=\mathrm{diag}(+1,-1,-1,-1,-1,-1,-1,-1):\qquad \Gamma_0^2=1_{16},\quad \Gamma_i^2=-1_{16}.
\end{equation*}
\textbf{Reason.} The task requires real matrices, and only this overall sign allows them; it is the sign pattern of the octonion basis, $\eta$ being simply the squares $e_0^2=+1$, $e_i^2=-1$ of the octonion units. With the opposite sign $\eta'=-\eta$ the relations $R_AR_B+R_BR_A=2\eta'_{AB}1_{16}$ define $Cl(7,1)\cong\mathbb H(8)$. If real 16×16 matrices $R_A$ satisfied them, then, because $Cl_8(\mathbb C)\cong\mathbb C(16)$ has only one irreducible module, $R_A=T(i\Gamma_A)T^{-1}$ for an invertible complex $T$ ($i\Gamma_A$ satisfy the same relations). $R_A=\overline{R_A}$ gives $Y\overline{(i\Gamma_A)}=(i\Gamma_A)Y$ with $Y=T^{-1}\overline T$, $Y\overline Y=1_{16}$; since $\overline{(i\Gamma_A)}=-i\Gamma_A$ this is $Y\Gamma_A=-\Gamma_AY$, so $Y\omega$ commutes with all $\Gamma_A$ and $Y=c\,\omega$ (Schur), whence $Y\overline Y=|c|^2\omega^2=-|c|^21_{16}\neq1_{16}$. Renaming, reordering or re-signing the indices does not change the signature of the relation (the numbers of generators squaring to $+1$ and to $-1$). The first stage proved this in 7+1.nb.

\textbf{The Dirac matrices with upper index, $\alpha$, $\beta$, $\omega$} (checks 1–5 of section 7.1 of 7+1.nb):
\begin{align*}
\Gamma^A&=\eta^{AA}\Gamma_A,\qquad (\Gamma^A)^\dagger=\Gamma^0\Gamma^A\Gamma^0,\qquad \beta=\Gamma^0=\begin{pmatrix}0&1_8\\1_8&0\end{pmatrix}, \\
\alpha_i&=\Gamma^0\Gamma^i=\begin{pmatrix}0&\tau^{O}_i\\-\tau^{O}_i&0\end{pmatrix},\qquad \alpha_i\alpha_j+\alpha_j\alpha_i=2\delta_{ij}1_{16},\qquad \alpha_i\beta+\beta\alpha_i=0,\qquad \beta^2=1_{16}, \\
\omega&=\begin{pmatrix}0&-1_8\\1_8&0\end{pmatrix},\qquad \omega^2=-1_{16},\qquad \omega\Gamma_A=-\Gamma_A\omega .
\end{align*}
\Needspace{8\baselineskip}\subsection*{3.2 The Dirac equation, the Klein–Gordon equation and the Hamiltonian}\addcontentsline{toc}{subsection}{3.2 The Dirac equation, the Klein–Gordon equation and the Hamiltonian}
\textbf{Decision.} $(i\Gamma^A\partial_A-m)\psi=0$. The factor $i$ is unavoidable: with real $\Gamma_A$, $(\Gamma^A\partial_A)^2=\tfrac12\lbrace\Gamma^A,\Gamma^B\rbrace\partial_A\partial_B=\Box$, so
\begin{equation*}
\big(i\Gamma^A\partial_A+m\big)\big(i\Gamma^A\partial_A-m\big)=-\big(\Box+m^2\big)1_{16},
\end{equation*}
which is the Klein–Gordon operator with the correct sign for a massive particle; without the $i$ the mass term would have the wrong sign ($(\Gamma^A\partial_A)^2-m^2=\Box-m^2$). Multiplying by $\Gamma^0$:
\begin{equation*}
\Gamma^0\big(i\Gamma^A\partial_A-m\big)\psi=\big(i\partial_0-h\big)\psi,\qquad h=-i\alpha_i\partial_i+\beta m,\qquad h^\dagger=h,
\end{equation*}
because $\alpha_i$, $\beta$ are real symmetric and $-i\partial_i$ is Hermitian. Conservation of $j^A=\bar\psi\Gamma^A\psi$, $\bar\psi=\psi^\dagger\Gamma^0$, follows from the adjoint equation $-i\partial_A\bar\psi\Gamma^A-m\bar\psi=0$: $\partial_Aj^A=(im\bar\psi)\psi+\bar\psi(-im\psi)=0$. 7+1.nb checks these identities for 16 arbitrary functions $\psi_a(x_0,\dots,x_7)$.

\Needspace{8\baselineskip}\subsection*{3.3 Real plane-wave spinors}\addcontentsline{toc}{subsection}{3.3 Real plane-wave spinors}
For $e^{-i(Ex_0-\vec p\cdot\vec x)}$, $h$ becomes the real symmetric matrix
\begin{align*}
h(\vec p)&=p^i\alpha_i+m\beta=\begin{pmatrix}0&M\\M^{\mathsf T}&0\end{pmatrix},\qquad M=m1_8+p^i\tau^{O}_i, \\
MM^{\mathsf T}&=m^21_8-(p^i\tau^{O}_i)^2=E^21_8,\qquad E=\sqrt{m^2+\vec p^{\,2}},
\end{align*}
using $(p^i\tau^{O}_i)^{\mathsf T}=-p^i\tau^{O}_i$ and $(p^i\tau^{O}_i)^2=-\vec p^{\,2}1_8$. \textbf{Derivation of the spinors.} With $u_s=\frac1{\sqrt E}(E\,e_s;\ M^{\mathsf T}e_s)$ and $v_s=\frac1{\sqrt E}(E\,e_s;\ -Me_s)$, and $h(-\vec p)=\left(\begin{smallmatrix}0&M^{\mathsf T}\\M&0\end{smallmatrix}\right)$:
\begin{align*}
h(\vec p)\,(E\,e_s;\ M^{\mathsf T}e_s)&=(MM^{\mathsf T}e_s;\ E\,M^{\mathsf T}e_s)=E\,(E\,e_s;\ M^{\mathsf T}e_s), \\
h(-\vec p)\,(E\,e_s;\ -Me_s)&=(-M^{\mathsf T}Me_s;\ E\,Me_s)=-E\,(E\,e_s;\ -Me_s), \\
u_s^{\mathsf T}u_r&=\frac{E^2\delta_{sr}+e_s^{\mathsf T}MM^{\mathsf T}e_r}{E}=2E\,\delta_{sr},\qquad v_s^{\mathsf T}v_r=2E\,\delta_{sr}.
\end{align*}
\textbf{Decision (normalization).} $u^{\mathsf T}u=2E$, the usual covariant normalization, so that $\sum_su_s\bar u_s=\not p+m$ and $\sum_sv_s\bar v_s=\not p-m$ with $\not p=E\Gamma_0+p^i\Gamma_i$. The example $\vec p=(2,2,4,0,0,0,0)$, $m=1$, $E=5$ was chosen because $E$ is an integer, so $\sqrt5\,u_s$ has integer entries. The real orthogonal matrix
\begin{equation*}
U=\frac1{\sqrt{2E}}\big(u_1(\vec p),\dots,u_8(\vec p),v_1(-\vec p),\dots,v_8(-\vec p)\big)\qquad\text{satisfies}\qquad U^{\mathsf T}h(\vec p)\,U=\mathrm{diag}(E1_8,-E1_8),
\end{equation*}
equivalently the completeness relation
\begin{equation*}
\sum_{s=1}^{8}u_s(\vec p)\,u_s(\vec p)^{\mathsf T}+\sum_{s=1}^{8}v_s(-\vec p)\,v_s(-\vec p)^{\mathsf T}=2E\,1_{16}.
\end{equation*}
\Needspace{8\baselineskip}\subsection*{3.4 Grassmann algebra, quantization and the canonical anticommutators}\addcontentsline{toc}{subsection}{3.4 Grassmann algebra, quantization and the canonical anticommutators}
\textbf{Decision.} The Grassmann algebra is implemented with ordered monomials \texttt{grass[\{k1, …, kr\}]}, with the product ($\pi$ the permutation that sorts $K\cup K'$) and the left derivative, as in Berezin:
\begin{align*}
\theta_K\theta_{K'}&=\mathrm{sign}(\pi)\,\theta_{\mathrm{sort}(K\cup K')}\qquad(\text{zero if }K\cap K'\neq\emptyset), \\
\frac{\partial^L(\theta_{k_1}\cdots\theta_{k_r})}{\partial\theta_{k_j}}&=(-1)^{j-1}\,\theta_{k_1}\cdots\widehat{\theta_{k_j}}\cdots\theta_{k_r} .
\end{align*}
The current-conservation check uses $\psi=\sum_{k=1}^3\theta_kU_k$, $\bar\psi=\sum_k\theta_{3+k}U_k^\dagger\Gamma^0$ with
\begin{align*}
U_1&=u_1(\vec p_1)e^{-i(5x_0-\vec p_1\cdot\vec x)},\qquad U_2=v_3(\vec p_3)e^{+i(3x_0-\vec p_3\cdot\vec x)},\qquad U_3=u_5(\vec0)e^{-ix_0}, \\
\vec p_1&=(2,2,4,0,0,0,0),\qquad \vec p_3=(0,0,0,2,2,0,0).
\end{align*}
The mode expansion and the anticommutators of the modes
\begin{align*}
\psi(x_0,\vec x)&=\sum_{\vec p}\sum_{s=1}^8\frac1{\sqrt{2EV}}\Big[b_s(\vec p)u_s(\vec p)e^{-i(Ex_0-\vec p\cdot\vec x)}+d_s^\dagger(\vec p)v_s(\vec p)e^{+i(Ex_0-\vec p\cdot\vec x)}\Big], \\
\lbrace b_s(\vec p),b_r^\dagger(\vec q)\rbrace&=\lbrace d_s(\vec p),d_r^\dagger(\vec q)\rbrace=\delta_{sr}\delta_{\vec p\vec q}\quad(\text{others }0)
\end{align*}
give, with the completeness relation of section 3.3 (the $v$ term contributes at $-\vec p$, which the symmetric sum over $\vec p$ relabels),
\begin{align*}
\lbrace\psi_a(x_0,\vec x),\psi_b^\dagger(x_0,\vec y)\rbrace&=\frac1V\sum_{\vec p}\frac{e^{i\vec p\cdot(\vec x-\vec y)}}{2E}\Big[\sum_su_su_s^{\mathsf T}+\sum_sv_s(-\vec p)v_s(-\vec p)^{\mathsf T}\Big]_{ab} \\
&=\delta_{ab}\frac1V\sum_{\vec p}e^{i\vec p\cdot(\vec x-\vec y)}=\delta_{ab}\,\delta^7(\vec x-\vec y).
\end{align*}
\textbf{Decision (test set).} The exact check uses five momenta closed under $\vec p\to-\vec p$, $\pm(2,2,4,0,0,0,0)$, $\pm(0,0,0,2,2,0,0)$, $\vec 0$ ($m=1$, $E=5,5,3,3,1$, all integers), with symbolic $\vec x$, $\vec y$ and box volume $V$. The normal-ordered Hamiltonian and charge follow from $U^{\mathsf T}hU=\mathrm{diag}(E,-E)$: $H=\sum E(b^\dagger b+d^\dagger d)$, $Q=\sum(b^\dagger b-d^\dagger d)$.

\Needspace{8\baselineskip}\subsection*{3.5 Charge conjugation}\addcontentsline{toc}{subsection}{3.5 Charge conjugation}
\textbf{Decision.} $\psi^c=\omega\psi^*$. Because the $\Gamma_A$ are real, conjugating the Dirac equation gives $(-i\Gamma^A\partial_A-m)\psi^*=0$, and $\omega$ is real and anticommutes with every $\Gamma_A$:
\begin{equation*}
\big(i\Gamma^A\partial_A-m\big)\omega\psi^*=\omega\big(-i\Gamma^A\partial_A-m\big)\psi^*=0,\qquad (\psi^c)^c=\omega\omega\psi=-\psi,\qquad (\not p+m)\,\omega\,u_s(\vec p)=0 .
\end{equation*}
Consequently $\psi^c=\psi$ has no non-zero solution: a massive real (Majorana) field needs 32 real components in 7+1 dimensions.

\Needspace{8\baselineskip}\subsection*{3.6 Lorentz covariance and the boost example}\addcontentsline{toc}{subsection}{3.6 Lorentz covariance and the boost example}
The spinor boost $\mathcal B(\beta_{\rm b})=e^{\beta_{\rm b}\Sigma_{01}}$ and its Lorentz matrix $L$ (section 2.1) satisfy $\mathcal B\Gamma_K\mathcal B^{-1}=\sum_BL_{BK}\Gamma_B$ and $L^{\mathsf T}\eta L=\eta$, with
\begin{equation*}
L=\begin{pmatrix}\cosh\beta_{\rm b}&-\sinh\beta_{\rm b}&0\\-\sinh\beta_{\rm b}&\cosh\beta_{\rm b}&0\\0&0&1_6\end{pmatrix}.
\end{equation*}
\textbf{Decision (example).} $\beta_{\rm b}=2\log2$, because then every number is rational: $\cosh\frac{\beta_{\rm b}}2=\frac54$, $\sinh\frac{\beta_{\rm b}}2=\frac34$, $\cosh\beta_{\rm b}=\frac{17}8$, $\sinh\beta_{\rm b}=\frac{15}8$. For $p=(5,2,2,4,0,0,0,0)$, $m=1$:
\begin{equation*}
p^{\prime}=Lp=\Big(\tfrac{17}{8}\cdot5-\tfrac{15}{8}\cdot2,\ -\tfrac{15}{8}\cdot5+\tfrac{17}{8}\cdot2,\ 2,4,0,0,0,0\Big)=\Big(\tfrac{55}{8},-\tfrac{41}{8},2,4,0,0,0,0\Big),\qquad p^{\prime}\cdot p^{\prime}=\tfrac{3025-1681}{64}-20=1,
\end{equation*}
and $(\not p\,'-m)\,\mathcal B\,u_s(\vec p)=0$ for $s=1,\dots,8$, with $\mathcal B=\frac54\,1_{16}+\frac34\,\Gamma_0\Gamma_1$.

\Needspace{8\baselineskip}\subsection*{3.7 Chirality and the two coupled 8-component equations}\addcontentsline{toc}{subsection}{3.7 Chirality and the two coupled 8-component equations}
\textbf{Result.} No real chirality matrix exists: the real 16×16 matrices $Z$ with $Z\Gamma_A=-\Gamma_AZ$ ($A=0,\dots,7$) form the space $\mathbb R\,\omega$ (7+1.nb solves the 256-unknown linear system, rank 255), and $(c\,\omega)^2=-c^21_{16}$ cannot be $1_{16}$. \textbf{Decision.} Chirality is the Hermitian matrix $\chi=i\omega$, unique up to sign; the sign is chosen so that $P_L=\frac12(1_{16}-\chi)$ projects on $\chi=-1$:
\begin{align*}
\chi&=\begin{pmatrix}0&-i1_8\\i1_8&0\end{pmatrix},\qquad \chi^2=1_{16},\qquad \chi\Gamma_A=-\Gamma_A\chi,\qquad \chi\Sigma_{AB}=\Sigma_{AB}\chi, \\
E_L&=\frac1{\sqrt2}\begin{pmatrix}1_8\\-i1_8\end{pmatrix},\quad E_R=\frac1{\sqrt2}\begin{pmatrix}-i1_8\\1_8\end{pmatrix},\qquad \chi E_L=-E_L,\quad \chi E_R=E_R,\quad E_LE_L^\dagger=P_L,\quad E_RE_R^\dagger=P_R .
\end{align*}
\textbf{Decision (chiral components).} $\phi_L=E_L^\dagger\psi=(\psi_\uparrow+i\psi_\downarrow)/\sqrt2$ and $\phi_R=E_R^\dagger\psi=(i\psi_\uparrow+\psi_\downarrow)/\sqrt2$. These isometries were chosen because they make the chiral blocks of the Dirac matrices as simple as possible:
\begin{align*}
E_L^\dagger\beta E_R&=E_R^\dagger\beta E_L=1_8,\qquad E_L^\dagger\beta E_L=E_R^\dagger\beta E_R=0, \\
E_L^\dagger\alpha_iE_L&=-\sigma_i,\qquad E_R^\dagger\alpha_iE_R=\sigma_i,\qquad E_L^\dagger\alpha_iE_R=0,\qquad \sigma_i=i\tau^{O}_i .
\end{align*}
\textbf{Derivation of the chiral equations.} Insert $\psi=E_L\phi_L+E_R\phi_R$ into $\Gamma^0(i\Gamma^A\partial_A-m)\psi=(i\partial_0+i\alpha_i\partial_i-\beta m)\psi$ and project with $E_L^\dagger$, $E_R^\dagger$; since $i\,(-\sigma_i)\partial_i=\tau^{O}_i\partial_i$:
\begin{align*}
E_L^\dagger\Gamma^0\big(i\Gamma^A\partial_A-m\big)\psi&=\big(i\partial_0+\tau^{O}_i\partial_i\big)\phi_L-m\,\phi_R, \\
E_R^\dagger\Gamma^0\big(i\Gamma^A\partial_A-m\big)\psi&=\big(i\partial_0-\tau^{O}_i\partial_i\big)\phi_R-m\,\phi_L .
\end{align*}
So the Dirac equation is equivalent to $(i\partial_0+\tau^{O}_i\partial_i)\phi_L=m\phi_R$ and $(i\partial_0-\tau^{O}_i\partial_i)\phi_R=m\phi_L$, coupled exactly when $m\neq0$, with only real 8×8 matrices besides the $i$ of $i\partial_0$. Each obeys the Klein–Gordon equation because $(i\partial_0-\tau^{O}_i\partial_i)(i\partial_0+\tau^{O}_j\partial_j)=-\Box1_8$. For a plane wave the equations read $(E+\sigma\cdot\vec p)\phi_L=m\phi_R$, $(E-\sigma\cdot\vec p)\phi_R=m\phi_L$, and $W=(E_L,E_R)$ gives $W^\dagger h(\vec p)W=\left(\begin{smallmatrix}-\sigma\cdot\vec p&m1_8\\m1_8&\sigma\cdot\vec p\end{smallmatrix}\right)$.

\Needspace{8\baselineskip}\subsection*{3.8 The self-coupled field}\addcontentsline{toc}{subsection}{3.8 The self-coupled field}
\textbf{Decision (interaction).} The field is coupled to itself by the four Lorentz-invariant contact (zero-range) interactions built from bilinears without derivatives that are used in relativistic point-coupling models: scalar, pseudoscalar, vector and axial vector,
\begin{align*}
\mathcal L&=\bar\psi\big(i\Gamma^A\partial_A-m\big)\psi-\bar\psi\big(\mathsf s+i\chi\,\mathsf p+\Gamma^A(\mathsf v_A+\chi\,\mathsf a_A)\big)\psi \\
&\quad-\frac{g_S}{2}(\bar\psi\psi)^2-\frac{g_P}{2}(\bar\psi\,i\chi\psi)^2-\frac{g_V}{2}(\bar\psi\Gamma^A\psi)(\bar\psi\Gamma_A\psi)-\frac{g_A}{2}(\bar\psi\Gamma^A\chi\psi)(\bar\psi\Gamma_A\chi\psi),
\end{align*}
with real external sources $\mathsf s,\mathsf p,\mathsf v_A,\mathsf a_A$ that define the densities. The tensor channel $(\bar\psi\Sigma_{AB}\psi)^2$ and the channels of higher rank are not included; they could be added in the same way, and by the Fierz identities of section 3.11 their exchange terms are again combinations of direct terms. The four densities are real because $\Gamma^0O$ is Hermitian for $O=1_{16},\ i\chi,\ \Gamma^A,\ \Gamma^A\chi$. \textbf{Units.} $\hbar=c=1$; $\psi$ has mass dimension $7/2$ and every $g$ mass dimension $-6$, so the theory is an effective one and needs either the no-sea approximation (section 3.10) or a momentum cutoff (section 3.14). \textbf{Signs.} The interaction has no time derivatives, so it enters the Hamiltonian density with the opposite sign:
\begin{align*}
\mathcal H&=\,:\psi^\dagger(-i\alpha_i\partial_i+\beta m)\psi:+:\bar\psi\big(\mathsf s+i\chi\mathsf p+\Gamma^A(\mathsf v_A+\chi\mathsf a_A)\big)\psi: \\
&\quad+\frac{g_S}2:\rho_S^2:+\frac{g_P}2:\rho_P^2:+\frac{g_V}2:j^Aj_A:+\frac{g_A}2:j_5^Aj_{5A}:,
\end{align*}
with $\rho_S=\bar\psi\psi$, $\rho_P=\bar\psi i\chi\psi$, $j^A=\bar\psi\Gamma^A\psi$, $j_5^A=\bar\psi\Gamma^A\chi\psi$; $g_S<0$ attracts.

\Needspace{8\baselineskip}\subsection*{3.9 The nonlinear field equation}\addcontentsline{toc}{subsection}{3.9 The nonlinear field equation}
The Euler–Lagrange equation is taken with the Grassmann LEFT derivative with respect to $\bar\psi_a$; the bilinears are Grassmann-even, so
\begin{align*}
&\big(i\Gamma^A\partial_A-m-\mathsf s-i\chi\mathsf p-\Gamma^A(\mathsf v_A+\chi\mathsf a_A)\big)\psi \\
&\qquad=g_S(\bar\psi\psi)\psi+g_P(\bar\psi i\chi\psi)\,i\chi\psi+g_V(\bar\psi\Gamma_A\psi)\Gamma^A\psi+g_A(\bar\psi\Gamma_A\chi\psi)\Gamma^A\chi\psi .
\end{align*}
\textbf{Decision (check).} 7+1.nb verifies $\partial^L\mathcal L_{\rm int}/\partial\bar\psi_a$ for all 16 components in a Grassmann algebra with the 32 generators $\psi_a=\theta_a$, $\bar\psi_a=\theta_{16+a}$, with symbolic couplings.

\Needspace{8\baselineskip}\subsection*{3.10 Hohenberg–Kohn theorem and Kohn–Sham equations}\addcontentsline{toc}{subsection}{3.10 Hohenberg–Kohn theorem and Kohn–Sham equations}
With $n=(\rho_S,\rho_P,j^A,j_5^A)$ (18 densities), $w=(\mathsf s,\mathsf p,\mathsf v_A,\mathsf a_A)$ and $H[w]=H_0+H_{\rm int}+\int w\cdot n\,d^7x$ in the sector of fixed charge $Q=N$, two sources $w\neq w'$ with different non-degenerate ground states $\Psi$, $\Psi'$ and the same densities would give
\begin{equation*}
E<\langle\Psi^{\prime}|H[w]|\Psi^{\prime}\rangle=E^{\prime}+\int(w-w^{\prime})\cdot n,\qquad E^{\prime}<E+\int(w^{\prime}-w)\cdot n\qquad\Rightarrow\qquad E+E^{\prime}<E+E^{\prime},
\end{equation*}
a contradiction; so $E[w]=\min_n\big(F[n]+\int w\cdot n\big)$ with $F[n]=\min_{\Psi\to n}\langle\Psi|H_0+H_{\rm int}|\Psi\rangle$. Writing $F=T_s+E_H+E_{xc}$ with
\begin{align*}
T_s&=\sum_{k=1}^N\int\varphi_k^\dagger(-i\alpha_i\partial_i+\beta m)\varphi_k,\qquad E_H=\int\Big[\frac{g_S}2\rho_S^2+\frac{g_P}2\rho_P^2+\frac{g_V}2j^Aj_A+\frac{g_A}2j_5^Aj_{5A}\Big],
\end{align*}
and varying $E=T_s+\int w\cdot n+E_H+E_{xc}$ with respect to $\varphi_k^\dagger$ under $\int\varphi_k^\dagger\varphi_l=\delta_{kl}$ (multipliers $\varepsilon_k$) gives $h_{\rm KS}\varphi_k=\varepsilon_k\varphi_k$ with
\begin{align*}
h_{\rm KS}&=\alpha_i(-i\partial_i+\mathcal V_i)+\beta(m+\mathcal S)+\mathcal V_0+\beta i\chi\mathcal P+\chi\mathcal A_0+\alpha_i\chi\mathcal A_i, \\
\mathcal S&=\mathsf s+g_S\rho_S+\frac{\delta E_{xc}}{\delta\rho_S},\quad \mathcal P=\mathsf p+g_P\rho_P+\frac{\delta E_{xc}}{\delta\rho_P},\quad \mathcal V_A=\mathsf v_A+g_Vj_A+\frac{\delta E_{xc}}{\delta j^A},\quad \mathcal A_A=\mathsf a_A+g_Aj_{5A}+\frac{\delta E_{xc}}{\delta j_5^A}.
\end{align*}
The variation uses $\delta\rho_S/\delta\varphi^\dagger=\beta\varphi$, $\delta\rho_P/\delta\varphi^\dagger=\beta i\chi\varphi$, $\delta j^A/\delta\varphi^\dagger=\Gamma^0\Gamma^A\varphi$, $\delta j_5^A/\delta\varphi^\dagger=\Gamma^0\Gamma^A\chi\varphi$, so the Hartree part is $v_H=\beta g_S\rho_S+\beta i\chi g_P\rho_P+\Gamma^0\Gamma^Ag_Vj_A+\Gamma^0\Gamma^A\chi g_Aj_{5A}$ (checked as Wirtinger derivatives, $\partial e_H/\partial\varphi_a^*=(v_H\varphi)_a$, with 32 independent variables). \textbf{Decision (no sea).} The occupied orbitals are the $N$ lowest positive-energy states; the total energy is
\begin{equation*}
E=\sum_k\varepsilon_k-\int\big[(\mathcal S-\mathsf s)\rho_S+(\mathcal P-\mathsf p)\rho_P+(\mathcal V_A-\mathsf v_A)j^A+(\mathcal A_A-\mathsf a_A)j_5^A\big]+E_H+E_{xc}.
\end{equation*}
\Needspace{8\baselineskip}\subsection*{3.11 Exchange (Wick) and the Fierz identities}\addcontentsline{toc}{subsection}{3.11 Exchange (Wick) and the Fierz identities}
For a Slater determinant of $N$ orbitals and contact vertices, Wick’s theorem with $\varrho_{xy}=\sum_k\varphi_{k,x}\varphi_{k,y}^*$ and $R=\varrho\Gamma^0$ gives
\begin{align*}
\langle:(\bar\psi O_1\psi)(\bar\psi O_2\psi):\rangle&=\mathrm{tr}(O_1R)\,\mathrm{tr}(O_2R)-\mathrm{tr}(O_1RO_2R),\qquad R=\sum_k\varphi_k\bar\varphi_k, \\
E_x&=-\frac{g_S}2\mathrm{tr}(RR)-\frac{g_P}2\mathrm{tr}(i\chi R\,i\chi R)-\frac{g_V}2\mathrm{tr}(\Gamma^AR\Gamma_AR)-\frac{g_A}2\mathrm{tr}(\Gamma^A\chi R\Gamma_A\chi R), \\
\Sigma_x&=-g_SR-g_P(i\chi)R(i\chi)-g_V\Gamma_AR\Gamma^A-g_A\Gamma_A\chi R\Gamma^A\chi,\qquad \delta E_x=\mathrm{tr}(\Sigma_x\delta R).
\end{align*}
\textbf{Decision (check).} 7+1.nb builds the 560-dimensional Fock space of 3 particles in 16 modes, with $c_a^\dagger c_b$ as explicit sparse matrices and $:B(M_1)B(M_2):=B(M_1)B(M_2)-B(M_1M_2)$, $B(M)=\sum M_{ab}c_a^\dagger c_b$, $M_i=\Gamma^0O_i$, and checks the formula for all 18 vertex pairs with 3 random orthonormal orbitals (\texttt{SeedRandom[71]}, 40 digits). The Fierz completeness $\sum_K(\Gamma_K)_{ab}(\Gamma_K^{-1})_{cd}=16\delta_{ad}\delta_{cb}$ gives the rearrangement
\begin{equation*}
\mathrm{tr}(O_1RO_2R)=\frac1{16}\sum_K\mathrm{tr}(\Gamma_K^{-1}R)\,\mathrm{tr}(O_1\Gamma_KO_2R),
\end{equation*}
checked with random complex integer matrices $R$, $\delta R$ (\texttt{SeedRandom[1771]}). \textbf{Uniform matter.} With $R=(\rho_S1_{16}+\rho\Gamma_0)/16$, using $\chi R\chi=(\rho_S-\rho\Gamma_0)/16$, $\Gamma^A\Gamma_A=8$ and $\Gamma^A\Gamma_0\Gamma_A=\Gamma_0+7\cdot(-\Gamma_0)=-6\Gamma_0$:
\begin{align*}
\mathrm{tr}(RR)&=\frac{\rho_S^2+\rho^2}{16},\qquad \mathrm{tr}(i\chi R\,i\chi R)=-\frac{\rho_S^2-\rho^2}{16}, \\
\mathrm{tr}(\Gamma^AR\Gamma_AR)&=\frac{8\rho_S^2-6\rho^2}{16},\qquad \mathrm{tr}(\Gamma^A\chi R\Gamma_A\chi R)=-\frac{8\rho_S^2+6\rho^2}{16}, \\
E_x&=-\frac{g_S}{32}(\rho_S^2+\rho^2)+\frac{g_P}{32}(\rho_S^2-\rho^2)-\frac{g_V}{32}(8\rho_S^2-6\rho^2)+\frac{g_A}{32}(8\rho_S^2+6\rho^2), \\
G_S&=g_S+\frac{-g_S+g_P-8g_V+8g_A}{16},\qquad G_V=g_V+\frac{-g_S-g_P+6g_V+6g_A}{16}, \\
\frac{g_S}2\rho_S^2+\frac{g_V}2\rho^2+E_x&=\frac{G_S}2\rho_S^2+\frac{G_V}2\rho^2
\end{align*}
(the pseudoscalar and axial Hartree terms vanish for uniform matter, where $\rho_P=j_5^A=0$). For $g_S=-10^5$, $g_V=10^4$, $g_P=g_A=0$: $G_S=-10^5+\frac{10^5-8\cdot10^4}{16}=-98750$ and $G_V=10^4+\frac{10^5+6\cdot10^4}{16}=20000$.

\Needspace{8\baselineskip}\subsection*{3.12 The chiral Kohn–Sham theory}\addcontentsline{toc}{subsection}{3.12 The chiral Kohn–Sham theory}
\textbf{Decision.} Use the variables $\kappa=\sum_k\varphi_{kL}^\dagger\varphi_{kR}$, $j_L^A=\sum_k\varphi_{kL}^\dagger\bar\sigma^A\varphi_{kL}$, $j_R^A=\sum_k\varphi_{kR}^\dagger\sigma^A\varphi_{kR}$ ($\sigma^A=(1_8,\sigma_i)$, $\bar\sigma^A=(1_8,-\sigma_i)$), the complex source $\varpi=\mathsf s+i\mathsf p$ and the complex mass $\mathcal M=m+\mathcal S+i\mathcal P$. From the chiral blocks of section 3.7:
\begin{align*}
\rho_S&=\kappa+\kappa^*,\quad \rho_P=i(\kappa-\kappa^*),\quad j^A=j_L^A+j_R^A,\quad j_5^A=j_R^A-j_L^A,\quad \mathsf s\rho_S+\mathsf p\rho_P=\varpi\kappa+\varpi^*\kappa^*, \\
e_H&=\frac{g_S-g_P}2(\kappa^2+\kappa^{*2})+(g_S+g_P)\kappa\kappa^*+\frac{g_V+g_A}2(j_L^Aj_{LA}+j_R^Aj_{RA})+(g_V-g_A)j_L^Aj_{RA}, \\
W^\dagger h_{\rm KS}(\vec p)W&=\begin{pmatrix}\mathcal V^L_0-\sigma_i(p^i+\mathcal V^L_i)&\mathcal M1_8\\ \mathcal M^*1_8&\mathcal V^R_0+\sigma_i(p^i+\mathcal V^R_i)\end{pmatrix},\qquad \mathcal V^{L,R}_A=\mathcal V_A\mp\mathcal A_A, \\
\mathcal M-m&=\varpi+\frac{\partial(e_H+e_{xc})}{\partial\kappa},\qquad \text{Hartree: }\ \mathcal M-m-\varpi=g_S(\kappa+\kappa^*)-g_P(\kappa-\kappa^*)=g_S\rho_S+ig_P\rho_P .
\end{align*}
\textbf{Notation decision.} The complex source is called $\varpi$, not $\mu$: $\mu$ is the chemical potential of section 3.13. \textbf{Chiral symmetry.} $e^{i\theta\chi}=\cos\theta\,1_{16}-\sin\theta\,\omega$ is real and multiplies $\kappa$ by $e^{2i\theta}$; $E_H$ changes by $\frac{g_S-g_P}2[(e^{4i\theta}-1)\kappa^2+(e^{-4i\theta}-1)\kappa^{*2}]$, so the interaction is chirally symmetric exactly for $g_P=g_S$. The check uses $\cos\theta=\frac35$, $\sin\theta=\frac45$ (rational).

\Needspace{8\baselineskip}\subsection*{3.13 Uniform matter: formulas, parameters and numerical precision}\addcontentsline{toc}{subsection}{3.13 Uniform matter: formulas, parameters and numerical precision}
Static uniform matter at rest, all 8 spin states filled to $k_F$, no sea, no sources: $\mathcal S=m^*-m$, $\mathcal V_0$ constant, $\mathcal P=\mathcal A_A=\mathcal V_i=0$. The area of the unit sphere $S^6\subset\mathbb R^7$ is $|S^6|=2\pi^{7/2}/\Gamma_{\rm E}(\frac72)=\frac{16\pi^3}{15}$ ($\Gamma_{\rm E}$ is Euler’s Gamma function, $\Gamma_{\rm E}(\frac72)=\frac{15}8\sqrt\pi$), so
\begin{align*}
8\int\frac{d^7k}{(2\pi)^7}f(|\vec k|)&=c_7\int_0^\infty k^6f(k)\,dk,\qquad c_7=\frac{8|S^6|}{(2\pi)^7}=\frac1{15\pi^4}, \\
\rho&=c_7\frac{k_F^7}7=\frac{k_F^7}{105\pi^4},\qquad \rho_S=c_7\int_0^{k_F}\frac{k^6m^*}{\sqrt{k^2+m^{*2}}}dk,\qquad \varepsilon_{\rm kin}=c_7\int_0^{k_F}k^6\sqrt{k^2+m^{*2}}\,dk, \\
m^*&=m+g_S\rho_S(m^*),\qquad \mathcal V_0=g_V\rho,\qquad \varepsilon_F=\mathcal V_0+\sqrt{k_F^2+m^{*2}}, \\
\varepsilon&=\varepsilon_{\rm kin}-\frac{(m^*-m)^2}{2g_S}+\frac{g_V}2\rho^2 .
\end{align*}
The integrals have the closed forms (obtained with \texttt{Integrate} in 7+1.nb, checked by differentiation with sympy), with $\mathrm{artanh}\,y=\frac12\log\frac{1+y}{1-y}$:
\begin{align*}
\int_0^{k_F}\frac{k^6m^*}{\sqrt{k^2+m^{*2}}}dk&=\frac{m^*}{48}\Big[k_F\sqrt{k_F^2+m^{*2}}\big(8k_F^4-10k_F^2m^{*2}+15m^{*4}\big) \\
&\qquad\qquad-15m^{*6}\,\mathrm{artanh}\frac{k_F}{\sqrt{k_F^2+m^{*2}}}\Big], \\
\int_0^{k_F}k^6\sqrt{k^2+m^{*2}}\,dk&=\frac1{384}\Big[k_F\sqrt{k_F^2+m^{*2}}\big(48k_F^6+8k_F^4m^{*2}-10k_F^2m^{*4}+15m^{*6}\big) \\
&\qquad\qquad-15m^{*8}\,\mathrm{artanh}\frac{k_F}{\sqrt{k_F^2+m^{*2}}}\Big].
\end{align*}
\textbf{Hugenholtz–Van Hove.} $\partial\varepsilon/\partial m^*=\rho_S-(m^*-m)/g_S=0$ is the gap equation (with $d\varepsilon_{\rm kin}/dm^*=\rho_S$), so $d\varepsilon/dk_F=c_7k_F^6(\sqrt{k_F^2+m^{*2}}+g_V\rho)$ and, with $d\rho/dk_F=c_7k_F^6$, the chemical potential is $\mu=d\varepsilon/d\rho=\varepsilon_F$; then $\frac{d}{dk_F}\frac{\varepsilon}{\rho}=\frac{c_7k_F^6}{\rho}\big(\mu-\frac\varepsilon\rho\big)$ vanishes at the minimum of $E/A=\varepsilon/\rho$, where $E/A=\varepsilon_F$.

\textbf{Decisions (parameters).} Units $m=1$. $g_S=-10^5$, $g_V=10^4$, $g_P=g_A=0$ (couplings of mass dimension $-6$): with these values the energy per particle has a minimum below $m$ at $k_F=0.7566\ldots$, so the example shows binding and saturation; the Hartree–Fock comparison is made at $k_F=\frac34$. \textbf{Decisions (precision).} The gap equation is solved by Brent’s method with 50-digit working precision (\texttt{wp = 50}, accuracy and precision goals $45$), the input $k_F$ padded to 60 digits; the saturation condition $\varepsilon/\rho=\varepsilon_F$ is solved by Brent’s method on $[0.7,0.8]$ with 40-digit working precision and goals 32. The checks require $|E/A-\varepsilon_F|<10^{-25}$ at saturation and compare $\mu$ from a symmetric difference quotient (step $h=10^{-12}$ in $k_F$) with $\varepsilon_F$ to $10^{-15}$. These settings were adopted after development runs in which \texttt{FindRoot} issued the message \texttt{FindRoot::precw} (input precision below the working precision) and its default goals (half the working precision) left residuals larger than the tolerances of the checks.

\Needspace{8\baselineskip}\subsection*{3.14 The chiral vacuum with the Dirac sea}\addcontentsline{toc}{subsection}{3.14 The chiral vacuum with the Dirac sea}
\textbf{Decision.} To show spontaneous chiral symmetry breaking, keep the Dirac sea with a momentum cutoff $|\vec p|<\Lambda$, take $g_P=g_S=g$, $\varpi=0$, $\mathcal V=\mathcal A=0$ and a uniform complex mass $\mathcal M$:
\begin{align*}
h_\chi(\vec p)&=\begin{pmatrix}-\sigma\cdot\vec p&\mathcal M1_8\\ \mathcal M^*1_8&\sigma\cdot\vec p\end{pmatrix},\qquad h_\chi^2=(\vec p^{\,2}+|\mathcal M|^2)1_{16},\qquad \Pi_-=\frac12\Big(1_{16}-\frac{h_\chi}{E}\Big),\quad E=\sqrt{\vec p^{\,2}+|\mathcal M|^2}, \\
\kappa_{\vec p}&=\mathrm{tr}\,\Pi_-^{RL}=-\frac{4\mathcal M^*}{E},\qquad \frac{|S^6|}{(2\pi)^7}=\frac1{120\pi^4},\qquad \kappa=-4\mathcal M^*J(|\mathcal M|), \\
J(M)&=\frac1{120\pi^4}\int_0^\Lambda\frac{p^6dp}{\sqrt{p^2+M^2}},\qquad K(M)=\frac1{120\pi^4}\int_0^\Lambda p^6\sqrt{p^2+M^2}\,dp, \\
\mathcal M&=m+2g\kappa^*=m-8gJ(|\mathcal M|)\mathcal M,\qquad \mathcal E(\mathcal M)=-8K+8J\,\mathrm{Re}\big[(\mathcal M-m)\mathcal M^*\big]+32g|\mathcal M|^2J^2 .
\end{align*}
Here $\Pi_-^{RL}$ is the lower-left 8×8 block of $\Pi_-$, equal to $-\mathcal M^*1_8/(2E)$, and $\mathcal E$ is the sum of the sea eigenvalues $-8K$ minus the double counting $2\mathrm{Re}[(\mathcal M-m)\kappa]$ plus $e_H=2g|\kappa|^2$. For $m=0$ a solution $\mathcal M\neq0$ needs $1+8gJ(|\mathcal M|)=0$, possible exactly when
\begin{equation*}
g<g_c=-\frac1{8J(0)}=-\frac{90\pi^4}{\Lambda^6},\qquad J(0)=\frac{\Lambda^6}{720\pi^4},\qquad K(0)=\frac{\Lambda^8}{960\pi^4}.
\end{equation*}
\textbf{Decisions (parameters).} $\Lambda=1$; $g=\frac65g_c$ (20\% beyond the critical coupling, so that $|\mathcal M|$ is clearly non-zero); a small explicit mass $m=\frac1{100}$ to show that it selects the real positive $\mathcal M$, $\mathcal E(\mathcal Me^{i\phi})-\mathcal E(\mathcal M)=8Jm|\mathcal M|(1-\cos\phi)>0$; the diagonalization check at $\vec p=(\frac3{10},-\frac15,\frac1{10},\frac14,0,-\frac3{20},\frac1{20})$ with $\mathcal M=|\mathcal M|e^{7i/10}$ (a generic momentum and phase). The gap equation is solved by Brent’s method with 40-digit working precision and goals 32.

\Needspace{8\baselineskip}\subsection*{3.15 Notation decisions}\addcontentsline{toc}{subsection}{3.15 Notation decisions}
\begin{itemize}
\item The coordinates $x_0,\dots,x_7$ and the metric $\eta$ of 7+1.nb are kept. Kohn–Sham potentials are calligraphic ($\mathcal S,\mathcal P,\mathcal V_A,\mathcal A_A$), sources sans-serif ($\mathsf s,\mathsf p,\mathsf v_A,\mathsf a_A$).
\item The Dirac matrix $\beta=\Gamma^0$ keeps its name; the rapidity of a boost is $\beta_{\rm b}$ (Mathematica \texttt{β} in \texttt{Sβ}), the spinor boost is $\mathcal B(\beta_{\rm b})$ (Mathematica \texttt{Sβ}) and its Lorentz matrix $L$ (Mathematica \texttt{Λboost}), so that $\mathcal S$ means only the scalar potential and $\Lambda$ only the cutoff.
\item The complex scalar source is $\varpi=\mathsf s+i\mathsf p$; $\mu$ is only the chemical potential.
\item The matrices of the two proofs are $Y$ (no real realization for $\eta'$) and $Z$ (no real chirality); $J$ and $K$ are only the sea integrals; Euler’s Gamma function is $\Gamma_{\rm E}$, to distinguish it from the Dirac matrices $\Gamma_A$.
\item Mathematica names: \texttt{Γ17[A]}, \texttt{βD}, \texttt{αD[i]}, \texttt{ω17}, \texttt{χ17}, \texttt{EL}, \texttt{ER}, \texttt{Wch}, \texttt{ϕL}, \texttt{ϕR}, \texttt{σW[i]}, \texttt{dirac16complex}, \texttt{hKS}, \texttt{hChiKS}, and so on; in the notebook text cells the same symbols are written in plain Unicode.
\end{itemize}
\Needspace{8\baselineskip}\subsection*{3.16 The signature: the author’s explanation of 27 October 2026}\addcontentsline{toc}{subsection}{3.16 The signature: the author’s explanation of 27 October 2026}
\textbf{Source.} On 27 October 2026 the author of the project explained that the xAct list $\{7,1,0\}$ declares seven space-like dimensions and one time-like dimension, while $\mathrm{diag}(+1,-1,\dots,-1)$ is simply the squares of the octonion basis, and asked that this be made clear in all documentation, examples and provenance records. \textbf{What it corrected.} The version of 7+1.nb pushed in the commits 52d8729 to dc2f178 of github.com/once-ere/7plus1 declared the xAct metric of section 12 with \texttt{DefMetric[\{1, 7, 0\}, …]}, reading the list as the signs of $\eta$; since then section 12 declares \texttt{DefMetric[\{7, 1, 0\}, g71[-a, -b], CD, …]}, exactly as the author’s notebook Einstein-Lovelock-7+1-Nash.nb does. \textbf{What was added.} Checks in 7+1.nb — section 4.7: $\eta$ is the squares of the octonion basis, $e_Ae_A=\eta_{AA}e_0$ and $\mathrm{Re}(e_Ae_B)=\eta_{AB}$; section 12.2: \texttt{SignatureOfMetric[g71]} $=\{7,1,0\}$ with \texttt{SignDetOfMetric[g71]} $=-1$, $\mathrm{Re}(e_Ae_B)=\eta_{AB}$, and $\mathrm{Re}(x\,x)=\eta_{AB}x^Ax^B$; section 17: the flat metric $\mathrm{diag}(+1,-1,-1,-1)$ of the world brane is the squares of the basis $e_0,e_1,e_2,e_3$ of the quaternion subalgebra — and the same statements in the Python programs \texttt{tools/\allowbreak{}ref71.py}, \texttt{tools/\allowbreak{}ref\_\allowbreak{}gravity.py} and \texttt{tools/\allowbreak{}ref\_\allowbreak{}sm.py}; the explanation itself in the text of 7+1.nb (sections 1.2, 1.4, 4.2, 4.7, 7, 11, 12.2, 13, 15.1, 17) and in every document, symbol file and provenance file of the repository. \textbf{What did not change.} No result. Compared output by output with the backup below, every output of 7+1.nb is the same, except the declared list itself, the new checks and the check counts (and, as in every run, the name of the private context and the measured times); in particular every component of every Lovelock tensor is unchanged. The components follow from the component matrix of the metric, and the sign of the determinant, \texttt{SignDetOfMetric}, is $-1$ for both lists. \textbf{Backup.} The evaluated notebook as it was before is kept in the workspace as \texttt{backups/\allowbreak{}7+1\_\allowbreak{}2026-10-27\_\allowbreak{}before\_\allowbreak{}signature.nb} (the folder backups/ is excluded from the repository by .gitignore), 1034401 bytes, SHA-256 \texttt{06094248\allowbreak{}8caf4f45\allowbreak{}120ffd37\allowbreak{}be0123e9\allowbreak{}6992a31c\allowbreak{}d7fdbde0\allowbreak{}9ee4dc4c\allowbreak{}ed4f5d3a}; \texttt{tools/\allowbreak{}verify\_\allowbreak{}7p1\_\allowbreak{}nb.py} compares every input, text and section cell of 7+1.nb with it and finds every cell present, each one unchanged or changed by this explanation, and it compares every output in the same way.

\textbf{The explanation, in full.} The 7+1 spacetime has one time-like dimension ($x_0$, the direction of the octonion unit $e_0=1$) and seven space-like dimensions ($x_1,x_2,x_3$ ordinary space, $x_4,\dots,x_7$ auxiliary space). Two different objects describe this, and neither may be replaced by the other. (i) The tensor package xAct, used in section 12 of 7+1.nb exactly as in the author’s notebook Einstein-Lovelock-7+1-Nash.nb, declares the metric with \texttt{DefMetric[\{7, 1, 0\}, g71[-a, -b], CD, …]}: the list $\{7,1,0\}$ declares seven space-like dimensions, one time-like dimension and no degenerate dimension — the numbers of entries $+1$ and $-1$ of $\eta'=\mathrm{diag}(-1,+1,\dots,+1)$ (\texttt{η71}), the same spacetime written in the sign convention of 4+4.nb (space-like $=+1$, time-like $=-1$). (ii) The matrix $\eta=-\eta'=\mathrm{diag}(+1,-1,\dots,-1)$ (\texttt{η17}) is simply the squares of the octonion basis: $e_0^2=+1$ for the unit $e_0=1$ (time) and $e_i^2=-1$ for the seven imaginary units $e_1,\dots,e_7$ (space); more generally it is the real part of the product of two basis units, and it gives the quadratic form of the octonions:
\begin{gather*}
\eta_{AB}=\mathrm{Re}(e_Ae_B),\qquad e_0^2=+1,\qquad e_i^2=-1\ \ (i=1,\dots,7),\\ \mathrm{Re}(x\,x)=x_0^2-x_1^2-\cdots-x_7^2=\eta_{AB}\,x^Ax^B\qquad (x=x^Ae_A\in\mathbb O).
\end{gather*}
$\eta$ gives the signs of the Clifford relation $\Gamma_A\Gamma_B+\Gamma_B\Gamma_A=2\,\eta_{AB}\,1_{16}$ of the real 16×16 Dirac matrices: time squares to $+1_{16}$ and each of the seven space directions to $-1_{16}$, exactly like the octonion units (in the notation $Cl(p,q)$, with $p$ generators squaring to $+1$ and $q$ to $-1$, the algebra $Cl(1,7)\cong\mathbb R(16)$). It also gives the signs of the metric components: $g_{00}>0$ along the time-like direction, negative squares along the space-like directions. Every component computed by 7+1.nb follows from the component matrix of the metric; the declared list fixes for xAct the sign of the determinant of the metric, \texttt{SignDetOfMetric[g71]} $=-1$, which is the same in both sign conventions, since $\det \eta=\det \eta'=-1$. 7+1.nb checks \texttt{SignatureOfMetric[g71]} $=\{7,1,0\}$, \texttt{SignDetOfMetric[g71]} $=-1$, $\mathrm{Re}(e_Ae_B)=\eta_{AB}$ and $\mathrm{Re}(x\,x)=\eta_{AB}x^Ax^B$.

\Needspace{12\baselineskip}\section*{4. Verification record}\addcontentsline{toc}{section}{4. Verification record}
\Needspace{8\baselineskip}\subsection*{4.1 The Mathematica run}\addcontentsline{toc}{subsection}{4.1 The Mathematica run}
{\footnotesize
\begin{longtable}{>{\raggedright\arraybackslash}p{\dimexpr 0.3617\linewidth-2\tabcolsep\relax}>{\raggedright\arraybackslash}p{\dimexpr 0.6383\linewidth-2\tabcolsep\relax}}
\toprule
\textbf{item} & \textbf{value} \\
\midrule\endfirsthead
\toprule
\textbf{item} & \textbf{value} \\
\midrule\endhead
\endfoot
\bottomrule\endlastfoot
program & \texttt{C:\textbackslash{}\allowbreak{}Program Files\textbackslash{}\allowbreak{}Wolfram Research\textbackslash{}\allowbreak{}Wolfram\textbackslash{}\allowbreak{}15.0.1\textbackslash{}\allowbreak{}WolframNB.exe} \\
\$Version & 15.0.1 for Microsoft Windows (64-bit) (July 2, 2026) \\
method & \texttt{tools/\allowbreak{}run\_\allowbreak{}eval.wl}: all cells evaluated in a fresh kernel (front-end kernel configuration "Eval71"; time limit 3600 s), then \texttt{NotebookSave}; requested through \texttt{tools/\allowbreak{}watch\_\allowbreak{}requests.wl} (job \texttt{eval}) \\
finished & Wed 28 Oct 2026 08:57:47 \\
evaluation time of the whole notebook & 283.883 s \\
private context of the notebook & \texttt{Notebook\$\$75\$191583} \\
input cells / output cells / message cells & 189 / 554 / 0 \\
checks in sections 7–9 / True & 74 / 74 \\
checks in the whole notebook / True & 479 / 479 \\
summary of section 10 & \texttt{Length[checks]} = \texttt{203}, \texttt{failedChecks} = \texttt{\{\}}, \texttt{allChecksPassed} = \texttt{True} \\
SHA-256 of the evaluated 7+1.nb (2127641 bytes) & \texttt{41d252f9\allowbreak{}852f210c\allowbreak{}228f45e1\allowbreak{}ac956d2d\allowbreak{}7ef26fae\allowbreak{}b0062a58\allowbreak{}80280e1a\allowbreak{}ae0effc7} \\*
\end{longtable}}
Before the final evaluation, every input cell of sections 7–9 was developed with \texttt{tools/\allowbreak{}eval\_\allowbreak{}cells\_\allowbreak{}log.wl}, which evaluates a generated notebook expression by expression in a fresh context, records the time, the messages and the result of every check, and stops at the first failure.

\Needspace{8\baselineskip}\subsection*{4.2 The checks of sections 7–9 of 7+1.nb and their results}\addcontentsline{toc}{subsection}{4.2 The checks of sections 7–9 of 7+1.nb and their results}
Each row is one cell \texttt{check["name", test]} of the notebook, in order; the test is the Wolfram Language expression in that cell, and the result is what Wolfram 15.0.1 returned.

\Needspace{6\baselineskip}\subsubsection*{7.1 The eight real Dirac matrices and the Dirac \texorpdfstring{\ensuremath{\alpha}}{α}, \texorpdfstring{\ensuremath{\beta}}{β}}
{\footnotesize
\begin{longtable}{>{\raggedright\arraybackslash}p{\dimexpr 0.0700\linewidth-2\tabcolsep\relax}>{\raggedright\arraybackslash}p{\dimexpr 0.8100\linewidth-2\tabcolsep\relax}>{\raggedright\arraybackslash}p{\dimexpr 0.1200\linewidth-2\tabcolsep\relax}}
\toprule
\textbf{\#} & \textbf{check} & \textbf{result} \\
\midrule\endfirsthead
\toprule
\textbf{\#} & \textbf{check} & \textbf{result} \\
\midrule\endhead
\endfoot
\bottomrule\endlastfoot
1 & \texttt{Dirac matrices: Γ17up[0] == Γ17[0], Γ17up[i] == -Γ17[i]; all REAL, every entry 0 or ±1} & True \\
2 & \texttt{Dirac adjoint: ConjugateTranspose[Γ17up[A]] == Γ17[0].Γ17up[A].Γ17[0] for A = 0..7} & True \\
3 & \texttt{βD == \{\{0,ID8\},\{ID8,0\}\} and αD[i] == \{\{0,τO[i]\},\{-τO[i],0\}\}: real and symmetric} & True \\
4 & \texttt{αD[i].αD[j] + αD[j].αD[i] == 2 δij ID16, αD[i].βD + βD.αD[i] == 0, βD.βD == ID16} & True \\
5 & \texttt{ω17 = Γ17[8] == \{\{0,-ID8\},\{ID8,0\}\} is real, ω17.ω17 == -ID16, and ω17 anticommutes with every Γ17[A]} & True \\*
\end{longtable}}
\Needspace{6\baselineskip}\subsubsection*{7.2 The field dirac16complex and its Dirac equation}
{\footnotesize
\begin{longtable}{>{\raggedright\arraybackslash}p{\dimexpr 0.0700\linewidth-2\tabcolsep\relax}>{\raggedright\arraybackslash}p{\dimexpr 0.8100\linewidth-2\tabcolsep\relax}>{\raggedright\arraybackslash}p{\dimexpr 0.1200\linewidth-2\tabcolsep\relax}}
\toprule
\textbf{\#} & \textbf{check} & \textbf{result} \\
\midrule\endfirsthead
\toprule
\textbf{\#} & \textbf{check} & \textbf{result} \\
\midrule\endhead
\endfoot
\bottomrule\endlastfoot
6 & \texttt{(I Γ17up[A] ∂\_\allowbreak{}A + m)(I Γ17up[A] ∂\_\allowbreak{}A - m) dirac16complex == -(∂0\textasciicircum{}2 - ∂1\textasciicircum{}2 - ... - ∂7\textasciicircum{}2 + m\textasciicircum{}2) dirac16complex (16 arbitrary functions)} & True \\*
\end{longtable}}
\Needspace{6\baselineskip}\subsubsection*{7.3 Hamiltonian form}
{\footnotesize
\begin{longtable}{>{\raggedright\arraybackslash}p{\dimexpr 0.0700\linewidth-2\tabcolsep\relax}>{\raggedright\arraybackslash}p{\dimexpr 0.8100\linewidth-2\tabcolsep\relax}>{\raggedright\arraybackslash}p{\dimexpr 0.1200\linewidth-2\tabcolsep\relax}}
\toprule
\textbf{\#} & \textbf{check} & \textbf{result} \\
\midrule\endfirsthead
\toprule
\textbf{\#} & \textbf{check} & \textbf{result} \\
\midrule\endhead
\endfoot
\bottomrule\endlastfoot
7 & \texttt{Γ17[0].(I Γ17up[A] ∂\_\allowbreak{}A - m) ψ == I ∂0 ψ - h ψ with h = -I αD[i] ∂\_\allowbreak{}i + βD m (16 arbitrary functions)} & True \\*
\end{longtable}}
\Needspace{6\baselineskip}\subsubsection*{7.4 Plane waves and the real spinors u\_s, v\_s}
{\footnotesize
\begin{longtable}{>{\raggedright\arraybackslash}p{\dimexpr 0.0700\linewidth-2\tabcolsep\relax}>{\raggedright\arraybackslash}p{\dimexpr 0.8100\linewidth-2\tabcolsep\relax}>{\raggedright\arraybackslash}p{\dimexpr 0.1200\linewidth-2\tabcolsep\relax}}
\toprule
\textbf{\#} & \textbf{check} & \textbf{result} \\
\midrule\endfirsthead
\toprule
\textbf{\#} & \textbf{check} & \textbf{result} \\
\midrule\endhead
\endfoot
\bottomrule\endlastfoot
8 & \texttt{h(p) == \{\{0, M\}, \{Transpose[M], 0\}\} with the REAL 8×8 matrix M = m ID8 + p\_\allowbreak{}i τO[i]} & True \\
9 & \texttt{M.Transpose[M] == Transpose[M].M == E\textasciicircum{}2 ID8 and h(p).h(p) == E\textasciicircum{}2 ID16, E\textasciicircum{}2 = m\textasciicircum{}2 + p\textasciicircum{}2} & True \\
10 & \texttt{h(p).u\_\allowbreak{}s == E u\_\allowbreak{}s and h(-p).v\_\allowbreak{}s == -E v\_\allowbreak{}s for s = 1..8} & True \\
11 & \texttt{(pslash - m).u\_\allowbreak{}s == 0 and (pslash + m).v\_\allowbreak{}s == 0, pslash = E Γ17[0] + p\_\allowbreak{}i Γ17[i]} & True \\
12 & \texttt{u\_\allowbreak{}s.u\_\allowbreak{}r == v\_\allowbreak{}s.v\_\allowbreak{}r == 2E δsr, and u\_\allowbreak{}s(p).v\_\allowbreak{}r(-p) == 0} & True \\
13 & \texttt{plane waves: u\_\allowbreak{}s Exp[-I(E x0 - p.x)] and v\_\allowbreak{}s Exp[+I(E x0 - p.x)] solve (I Γ17up[A] ∂\_\allowbreak{}A - m) ψ == 0} & True \\
14 & \texttt{example p = (2,2,4,0,0,0,0), m = 1: E = 5 and the eight u\_\allowbreak{}s are real (entries in Sqrt[5] ℚ) and solve h(p).u\_\allowbreak{}s == 5 u\_\allowbreak{}s} & True \\*
\end{longtable}}
\Needspace{6\baselineskip}\subsubsection*{7.5 Spin sums, completeness, and the diagonalization of h}
{\footnotesize
\begin{longtable}{>{\raggedright\arraybackslash}p{\dimexpr 0.0700\linewidth-2\tabcolsep\relax}>{\raggedright\arraybackslash}p{\dimexpr 0.8100\linewidth-2\tabcolsep\relax}>{\raggedright\arraybackslash}p{\dimexpr 0.1200\linewidth-2\tabcolsep\relax}}
\toprule
\textbf{\#} & \textbf{check} & \textbf{result} \\
\midrule\endfirsthead
\toprule
\textbf{\#} & \textbf{check} & \textbf{result} \\
\midrule\endhead
\endfoot
\bottomrule\endlastfoot
15 & \texttt{spin sums: Σ\_\allowbreak{}s u\_\allowbreak{}s ubar\_\allowbreak{}s == pslash + m and Σ\_\allowbreak{}s v\_\allowbreak{}s vbar\_\allowbreak{}s == pslash - m} & True \\
16 & \texttt{completeness: Σ\_\allowbreak{}s u\_\allowbreak{}s(p) u\_\allowbreak{}s(p)\textasciicircum{}T + Σ\_\allowbreak{}s v\_\allowbreak{}s(-p) v\_\allowbreak{}s(-p)\textasciicircum{}T == 2E ID16} & True \\
17 & \texttt{U = (u\_\allowbreak{}1..u\_\allowbreak{}8, v\_\allowbreak{}1(-p)..v\_\allowbreak{}8(-p))/\allowbreak{}Sqrt[2E] is orthogonal and Transpose[U].h(p).U == diag(E ID8, -E ID8)} & True \\*
\end{longtable}}
\Needspace{6\baselineskip}\subsubsection*{7.6 Grassmann algebra: anticommuting components, Lagrangian, conserved current}
{\footnotesize
\begin{longtable}{>{\raggedright\arraybackslash}p{\dimexpr 0.0700\linewidth-2\tabcolsep\relax}>{\raggedright\arraybackslash}p{\dimexpr 0.8100\linewidth-2\tabcolsep\relax}>{\raggedright\arraybackslash}p{\dimexpr 0.1200\linewidth-2\tabcolsep\relax}}
\toprule
\textbf{\#} & \textbf{check} & \textbf{result} \\
\midrule\endfirsthead
\toprule
\textbf{\#} & \textbf{check} & \textbf{result} \\
\midrule\endhead
\endfoot
\bottomrule\endlastfoot
18 & \texttt{Grassmann: θ1 θ2 == -θ2 θ1, θ1 θ1 == 0, associativity, and ∂/\allowbreak{}∂θ2 (θ1 θ2 θ3) == -θ1 θ3} & True \\
19 & \texttt{the three plane waves U1, U2, U3 solve the Dirac equation with m = 1} & True \\
20 & \texttt{Grassmann field: ψ\_\allowbreak{}a ψ\_\allowbreak{}b + ψ\_\allowbreak{}b ψ\_\allowbreak{}a == 0 for all a, b (in particular ψ\_\allowbreak{}a\textasciicircum{}2 == 0)} & True \\
21 & \texttt{current conservation ∂\_\allowbreak{}A (ψbar Γ17up[A] ψ) == 0 for the Grassmann-valued solution ψ = θ1 U1 + θ2 U2 + θ3 U3} & True \\*
\end{longtable}}
\Needspace{6\baselineskip}\subsubsection*{7.7 Quantization: mode expansion and canonical anticommutators}
{\footnotesize
\begin{longtable}{>{\raggedright\arraybackslash}p{\dimexpr 0.0700\linewidth-2\tabcolsep\relax}>{\raggedright\arraybackslash}p{\dimexpr 0.8100\linewidth-2\tabcolsep\relax}>{\raggedright\arraybackslash}p{\dimexpr 0.1200\linewidth-2\tabcolsep\relax}}
\toprule
\textbf{\#} & \textbf{check} & \textbf{result} \\
\midrule\endfirsthead
\toprule
\textbf{\#} & \textbf{check} & \textbf{result} \\
\midrule\endhead
\endfoot
\bottomrule\endlastfoot
22 & \texttt{canonical anticommutator from the mode algebra: \{ψ\_\allowbreak{}a(x0,x), ψ\_\allowbreak{}b(x0,y)\textasciicircum{}†\} == δab (1/\allowbreak{}V) Σ\_\allowbreak{}p Exp[I p.(x - y)] (five momenta, symbolic x, y)} & True \\*
\end{longtable}}
\Needspace{6\baselineskip}\subsubsection*{7.8 Charge conjugation}
{\footnotesize
\begin{longtable}{>{\raggedright\arraybackslash}p{\dimexpr 0.0700\linewidth-2\tabcolsep\relax}>{\raggedright\arraybackslash}p{\dimexpr 0.8100\linewidth-2\tabcolsep\relax}>{\raggedright\arraybackslash}p{\dimexpr 0.1200\linewidth-2\tabcolsep\relax}}
\toprule
\textbf{\#} & \textbf{check} & \textbf{result} \\
\midrule\endfirsthead
\toprule
\textbf{\#} & \textbf{check} & \textbf{result} \\
\midrule\endhead
\endfoot
\bottomrule\endlastfoot
23 & \texttt{ψ -> ψc = ω17.ψ\textasciicircum{}*: (I Γ17up[A] ∂\_\allowbreak{}A - m)(ω17.ψ\textasciicircum{}*) == ω17.((I Γ17up[A] ∂\_\allowbreak{}A - m) ψ)\textasciicircum{}* for arbitrary ψ = ψre + I ψim} & True \\
24 & \texttt{(pslash + m).ω17.u\_\allowbreak{}s == 0 and ω17.u\_\allowbreak{}s == Σ\_\allowbreak{}r v\_\allowbreak{}r (-Transpose[M]/\allowbreak{}E)\_\allowbreak{}rs; ω17.ω17 == -ID16 (no Majorana spinors)} & True \\*
\end{longtable}}
\Needspace{6\baselineskip}\subsubsection*{7.9 Lorentz covariance}
{\footnotesize
\begin{longtable}{>{\raggedright\arraybackslash}p{\dimexpr 0.0700\linewidth-2\tabcolsep\relax}>{\raggedright\arraybackslash}p{\dimexpr 0.8100\linewidth-2\tabcolsep\relax}>{\raggedright\arraybackslash}p{\dimexpr 0.1200\linewidth-2\tabcolsep\relax}}
\toprule
\textbf{\#} & \textbf{check} & \textbf{result} \\
\midrule\endfirsthead
\toprule
\textbf{\#} & \textbf{check} & \textbf{result} \\
\midrule\endhead
\endfoot
\bottomrule\endlastfoot
25 & \texttt{boost at β = 2 log 2: S == 5/\allowbreak{}4 ID16 + 3/\allowbreak{}4 Γ17[0].Γ17[1], p' = Λ.p == (55/\allowbreak{}8, -41/\allowbreak{}8, 2, 4, 0, 0, 0, 0) is on shell, (pslash' - m).S.u\_\allowbreak{}s(p) == 0} & True \\*
\end{longtable}}
\Needspace{6\baselineskip}\subsubsection*{7.10 Chirality: no real chirality matrix, complex chiral projections}
{\footnotesize
\begin{longtable}{>{\raggedright\arraybackslash}p{\dimexpr 0.0700\linewidth-2\tabcolsep\relax}>{\raggedright\arraybackslash}p{\dimexpr 0.8100\linewidth-2\tabcolsep\relax}>{\raggedright\arraybackslash}p{\dimexpr 0.1200\linewidth-2\tabcolsep\relax}}
\toprule
\textbf{\#} & \textbf{check} & \textbf{result} \\
\midrule\endfirsthead
\toprule
\textbf{\#} & \textbf{check} & \textbf{result} \\
\midrule\endhead
\endfoot
\bottomrule\endlastfoot
26 & \texttt{χ17 = I ω17: Hermitian, χ17.χ17 == ID16, anticommutes with every Γ17[A], commutes with every Σ17[A,B]} & True \\
27 & \texttt{PL17, PR17 are complex (not real); Γ17[A].PL17 == PR17.Γ17[A]; ConjugateTranspose[PL17].Γ17[0] == Γ17[0].PR17} & True \\
28 & \texttt{EL, ER: χ17.EL == -EL, χ17.ER == ER, EL\textasciicircum{}†.EL == ER\textasciicircum{}†.ER == ID8, EL\textasciicircum{}†.ER == 0, EL.EL\textasciicircum{}† == PL17, ER.ER\textasciicircum{}† == PR17, W unitary} & True \\
29 & \texttt{φL == (ψup + I ψdown)/\allowbreak{}Sqrt[2], φR == (I ψup + ψdown)/\allowbreak{}Sqrt[2], PL17.ψ == EL.φL, PR17.ψ == ER.φR, ψ == EL.φL + ER.φR} & True \\*
\end{longtable}}
\Needspace{6\baselineskip}\subsubsection*{7.11 The two coupled 8-component equations}
{\footnotesize
\begin{longtable}{>{\raggedright\arraybackslash}p{\dimexpr 0.0700\linewidth-2\tabcolsep\relax}>{\raggedright\arraybackslash}p{\dimexpr 0.8100\linewidth-2\tabcolsep\relax}>{\raggedright\arraybackslash}p{\dimexpr 0.1200\linewidth-2\tabcolsep\relax}}
\toprule
\textbf{\#} & \textbf{check} & \textbf{result} \\
\midrule\endfirsthead
\toprule
\textbf{\#} & \textbf{check} & \textbf{result} \\
\midrule\endhead
\endfoot
\bottomrule\endlastfoot
30 & \texttt{σW[i] = I τO[i]: Hermitian, σW[i].σW[j] + σW[j].σW[i] == 2 δij ID8} & True \\
31 & \texttt{chiral (Weyl) form: W\textasciicircum{}†.h(p).W == \{\{-σW.p, m ID8\}, \{m ID8, σW.p\}\} (symbolic p)} & True \\
32 & \texttt{EL\textasciicircum{}†.Γ17[0].(I Γ17up[A] ∂\_\allowbreak{}A - m) ψ == (I ∂0 + τO[i] ∂\_\allowbreak{}i) φL - m φR (16 arbitrary functions)} & True \\
33 & \texttt{ER\textasciicircum{}†.Γ17[0].(I Γ17up[A] ∂\_\allowbreak{}A - m) ψ == (I ∂0 - τO[i] ∂\_\allowbreak{}i) φR - m φL (16 arbitrary functions)} & True \\
34 & \texttt{Klein-Gordon for each 8-component field: (I ∂0 - τO[i] ∂\_\allowbreak{}i)(I ∂0 + τO[j] ∂\_\allowbreak{}j) f == -(∂0\textasciicircum{}2 - ∂1\textasciicircum{}2 - ... - ∂7\textasciicircum{}2) f (8 arbitrary functions)} & True \\
35 & \texttt{plane-wave example p = (2,2,4,0,0,0,0), m = 1, E = 5: (E + σW.p) φL == m φR and (E - σW.p) φR == m φL for all eight u\_\allowbreak{}s} & True \\*
\end{longtable}}
\Needspace{6\baselineskip}\subsubsection*{8.1 The self-coupled field: Lagrangian, densities, Hamiltonian}
{\footnotesize
\begin{longtable}{>{\raggedright\arraybackslash}p{\dimexpr 0.0700\linewidth-2\tabcolsep\relax}>{\raggedright\arraybackslash}p{\dimexpr 0.8100\linewidth-2\tabcolsep\relax}>{\raggedright\arraybackslash}p{\dimexpr 0.1200\linewidth-2\tabcolsep\relax}}
\toprule
\textbf{\#} & \textbf{check} & \textbf{result} \\
\midrule\endfirsthead
\toprule
\textbf{\#} & \textbf{check} & \textbf{result} \\
\midrule\endhead
\endfoot
\bottomrule\endlastfoot
36 & \texttt{the densities ρS, ρP, j\textasciicircum{}A, j5\textasciicircum{}A are real: Γ17[0].op is Hermitian for op = ID16, I χ17, Γ17up[A], Γ17up[A].χ17} & True \\
37 & \texttt{18 vertex pairs; lower vertex == η17[[A+1,A+1]] upper vertex in the V and A channels} & True \\*
\end{longtable}}
\Needspace{6\baselineskip}\subsubsection*{8.2 The nonlinear field equation (Grassmann left derivatives)}
{\footnotesize
\begin{longtable}{>{\raggedright\arraybackslash}p{\dimexpr 0.0700\linewidth-2\tabcolsep\relax}>{\raggedright\arraybackslash}p{\dimexpr 0.8100\linewidth-2\tabcolsep\relax}>{\raggedright\arraybackslash}p{\dimexpr 0.1200\linewidth-2\tabcolsep\relax}}
\toprule
\textbf{\#} & \textbf{check} & \textbf{result} \\
\midrule\endfirsthead
\toprule
\textbf{\#} & \textbf{check} & \textbf{result} \\
\midrule\endhead
\endfoot
\bottomrule\endlastfoot
38 & \texttt{nonlinear field equation: left derivative ∂L\_\allowbreak{}int/\allowbreak{}∂ψbar\_\allowbreak{}a == -[gS (ψbar ψ) ψ + gP (ψbar Iχ17 ψ) Iχ17 ψ + gV (ψbar Γ17[A] ψ) Γ17up[A] ψ + gA (ψbar Γ17[A]χ17 ψ) Γ17up[A]χ17 ψ]\_\allowbreak{}a (32 Grassmann generators)} & True \\*
\end{longtable}}
\Needspace{6\baselineskip}\subsubsection*{8.4 The 16-component Kohn–Sham equations}
{\footnotesize
\begin{longtable}{>{\raggedright\arraybackslash}p{\dimexpr 0.0700\linewidth-2\tabcolsep\relax}>{\raggedright\arraybackslash}p{\dimexpr 0.8100\linewidth-2\tabcolsep\relax}>{\raggedright\arraybackslash}p{\dimexpr 0.1200\linewidth-2\tabcolsep\relax}}
\toprule
\textbf{\#} & \textbf{check} & \textbf{result} \\
\midrule\endfirsthead
\toprule
\textbf{\#} & \textbf{check} & \textbf{result} \\
\midrule\endhead
\endfoot
\bottomrule\endlastfoot
39 & \texttt{h\_\allowbreak{}KS(p) is Hermitian for real S, P, V\_\allowbreak{}A, A\_\allowbreak{}A (symbolic momentum and potentials)} & True \\
40 & \texttt{plane wave: Γ17[0].(pslash(ε,p) - m - S - I χ17 P - Γ17up[A](V\_\allowbreak{}A + χ17 A\_\allowbreak{}A)) == ε ID16 - h\_\allowbreak{}KS(p)} & True \\
41 & \texttt{Hartree potentials: ∂E\_\allowbreak{}H/\allowbreak{}∂φ\_\allowbreak{}a\textasciicircum{}* == (v\_\allowbreak{}H.φ)\_\allowbreak{}a, v\_\allowbreak{}H = h\_\allowbreak{}KS - h with S = gS ρS, P = gP ρP, V\_\allowbreak{}A = gV j\_\allowbreak{}A, A\_\allowbreak{}A = gA j5\_\allowbreak{}A (32 variables)} & True \\*
\end{longtable}}
\Needspace{6\baselineskip}\subsubsection*{8.5 Exchange: Wick’s theorem for contact couplings}
{\footnotesize
\begin{longtable}{>{\raggedright\arraybackslash}p{\dimexpr 0.0700\linewidth-2\tabcolsep\relax}>{\raggedright\arraybackslash}p{\dimexpr 0.8100\linewidth-2\tabcolsep\relax}>{\raggedright\arraybackslash}p{\dimexpr 0.1200\linewidth-2\tabcolsep\relax}}
\toprule
\textbf{\#} & \textbf{check} & \textbf{result} \\
\midrule\endfirsthead
\toprule
\textbf{\#} & \textbf{check} & \textbf{result} \\
\midrule\endhead
\endfoot
\bottomrule\endlastfoot
42 & \texttt{Wick's theorem in the 560-dimensional Fock space: <Φ| :(ψbar o1 ψ)(ψbar o2 ψ): |Φ> == tr(o1 R) tr(o2 R) - tr(o1 R o2 R) for all 18 vertex pairs} & True \\*
\end{longtable}}
\Needspace{6\baselineskip}\subsubsection*{8.6 The exchange self-energy and the Fierz identities}
{\footnotesize
\begin{longtable}{>{\raggedright\arraybackslash}p{\dimexpr 0.0700\linewidth-2\tabcolsep\relax}>{\raggedright\arraybackslash}p{\dimexpr 0.8100\linewidth-2\tabcolsep\relax}>{\raggedright\arraybackslash}p{\dimexpr 0.1200\linewidth-2\tabcolsep\relax}}
\toprule
\textbf{\#} & \textbf{check} & \textbf{result} \\
\midrule\endfirsthead
\toprule
\textbf{\#} & \textbf{check} & \textbf{result} \\
\midrule\endhead
\endfoot
\bottomrule\endlastfoot
43 & \texttt{exchange self-energy: d/\allowbreak{}dt E\_\allowbreak{}x(R + t dR) at t = 0 == tr(Σ\_\allowbreak{}x.dR) and E\_\allowbreak{}x == tr(Σ\_\allowbreak{}x.R)/\allowbreak{}2 (random complex R, dR; symbolic couplings)} & True \\
44 & \texttt{the 256 products Γ\_\allowbreak{}K are orthogonal matrices with tr(Inverse[Γ\_\allowbreak{}K].Γ\_\allowbreak{}L) == 16 δKL} & True \\
45 & \texttt{Fierz completeness: Σ\_\allowbreak{}K (Γ\_\allowbreak{}K)\_\allowbreak{}ab (Inverse[Γ\_\allowbreak{}K])\_\allowbreak{}cd == 16 δad δcb (as the 256×256 matrix identity Transpose[F].F == 16 ID)} & True \\
46 & \texttt{Fierz rearrangement of every exchange term: tr(o1 R o2 R) == (1/\allowbreak{}16) Σ\_\allowbreak{}K tr(Inverse[Γ\_\allowbreak{}K].R) tr(o1.Γ\_\allowbreak{}K.o2.R) for all 18 vertex pairs (random complex R)} & True \\*
\end{longtable}}
\Needspace{6\baselineskip}\subsubsection*{8.7 Example: uniform matter (no sea), Hartree and Hartree–Fock}
{\footnotesize
\begin{longtable}{>{\raggedright\arraybackslash}p{\dimexpr 0.0700\linewidth-2\tabcolsep\relax}>{\raggedright\arraybackslash}p{\dimexpr 0.8100\linewidth-2\tabcolsep\relax}>{\raggedright\arraybackslash}p{\dimexpr 0.1200\linewidth-2\tabcolsep\relax}}
\toprule
\textbf{\#} & \textbf{check} & \textbf{result} \\
\midrule\endfirsthead
\toprule
\textbf{\#} & \textbf{check} & \textbf{result} \\
\midrule\endhead
\endfoot
\bottomrule\endlastfoot
47 & \texttt{uniform matter: tr R == ρS, tr(Γ17up[0].R) == ρ, and the densities ρP, j\textasciicircum{}i, j5\textasciicircum{}A vanish} & True \\
48 & \texttt{the occupied plane waves give this R: (1/\allowbreak{}2) Σ\_\allowbreak{}s [u\_\allowbreak{}s(p) ubar\_\allowbreak{}s(p) + u\_\allowbreak{}s(-p) ubar\_\allowbreak{}s(-p)]/\allowbreak{}(2E) == (E Γ17[0] + m ID16)/\allowbreak{}(2E)} & True \\
49 & \texttt{uniform exchange energy: E\_\allowbreak{}x(R) == -gS/\allowbreak{}32 (ρS\textasciicircum{}2 + ρ\textasciicircum{}2) + gP/\allowbreak{}32 (ρS\textasciicircum{}2 - ρ\textasciicircum{}2) - gV/\allowbreak{}32 (8ρS\textasciicircum{}2 - 6ρ\textasciicircum{}2) + gA/\allowbreak{}32 (8ρS\textasciicircum{}2 + 6ρ\textasciicircum{}2)} & True \\
50 & \texttt{uniform exchange self-energy: Σ\_\allowbreak{}x == S\_\allowbreak{}x ID16 + V\_\allowbreak{}x Γ17[0] with S\_\allowbreak{}x == ∂E\_\allowbreak{}x/\allowbreak{}∂ρS and V\_\allowbreak{}x == ∂E\_\allowbreak{}x/\allowbreak{}∂ρ} & True \\
51 & \texttt{uniform Hartree-Fock == Hartree with the Fierz-transformed couplings: gS/\allowbreak{}2 ρS\textasciicircum{}2 + gV/\allowbreak{}2 ρ\textasciicircum{}2 + E\_\allowbreak{}x == G\_\allowbreak{}S/\allowbreak{}2 ρS\textasciicircum{}2 + G\_\allowbreak{}V/\allowbreak{}2 ρ\textasciicircum{}2 (gP, gA arbitrary)} & True \\
52 & \texttt{c7 == 8 |S\textasciicircum{}6|/\allowbreak{}(2 Pi)\textasciicircum{}7 == 1/\allowbreak{}(15 Pi\textasciicircum{}4), |S\textasciicircum{}6| == 2 Pi\textasciicircum{}(7/\allowbreak{}2)/\allowbreak{}Gamma[7/\allowbreak{}2] == 16 Pi\textasciicircum{}3/\allowbreak{}15, and ρ == c7 kF\textasciicircum{}7/\allowbreak{}7 == kF\textasciicircum{}7/\allowbreak{}(105 Pi\textasciicircum{}4)} & True \\
53 & \texttt{closed forms: d/\allowbreak{}dkF of the integrals gives the integrands, both vanish at kF = 0, and dε\_\allowbreak{}kin/\allowbreak{}dm* == ρS} & True \\
54 & \texttt{Hartree gap equation m* == m + gS ρS(m*) solved to 45 digits; m* is a stationary point of ε(m*) at fixed kF} & True \\
55 & \texttt{total energy == sum of KS eigenvalues - E\_\allowbreak{}H (double counting) == ε\_\allowbreak{}kin - (m* - m) ρS + E\_\allowbreak{}H} & True \\
56 & \texttt{Hugenholtz-Van Hove at saturation: E/\allowbreak{}A == ε\_\allowbreak{}F == dε/\allowbreak{}dρ; E/\allowbreak{}A - m < 0 and E/\allowbreak{}A at 0.95 kF0 and 1.05 kF0 is higher} & True \\
57 & \texttt{Hugenholtz-Van Hove away from saturation: μ == dε/\allowbreak{}dρ == ε\_\allowbreak{}F at kF = 1} & True \\
58 & \texttt{Hartree-Fock at kF = 3/\allowbreak{}4 (gP = gA = 0): G\_\allowbreak{}S == -98750, G\_\allowbreak{}V == 20000; m* - m == gS ρS + S\_\allowbreak{}x, V\_\allowbreak{}0 == gV ρ + V\_\allowbreak{}x; μ == ε\_\allowbreak{}F} & True \\*
\end{longtable}}
\Needspace{6\baselineskip}\subsubsection*{9.1 Chiral densities}
{\footnotesize
\begin{longtable}{>{\raggedright\arraybackslash}p{\dimexpr 0.0700\linewidth-2\tabcolsep\relax}>{\raggedright\arraybackslash}p{\dimexpr 0.8100\linewidth-2\tabcolsep\relax}>{\raggedright\arraybackslash}p{\dimexpr 0.1200\linewidth-2\tabcolsep\relax}}
\toprule
\textbf{\#} & \textbf{check} & \textbf{result} \\
\midrule\endfirsthead
\toprule
\textbf{\#} & \textbf{check} & \textbf{result} \\
\midrule\endhead
\endfoot
\bottomrule\endlastfoot
59 & \texttt{chiral densities: ρS == κ + κ\textasciicircum{}*, ρP == I(κ - κ\textasciicircum{}*), j\textasciicircum{}A == jL\textasciicircum{}A + jR\textasciicircum{}A, j5\textasciicircum{}A == jR\textasciicircum{}A - jL\textasciicircum{}A (symbolic chiral orbitals)} & True \\
60 & \texttt{the chiral blocks behind these densities: EL\textasciicircum{}†.βD.ER == ID8, EL\textasciicircum{}†.βD.EL == 0, EL\textasciicircum{}†.αD[i].EL == -σW[i], ER\textasciicircum{}†.αD[i].ER == σW[i], EL\textasciicircum{}†.αD[i].ER == 0} & True \\*
\end{longtable}}
\Needspace{6\baselineskip}\subsubsection*{9.2 The Hartree energy in chiral variables}
{\footnotesize
\begin{longtable}{>{\raggedright\arraybackslash}p{\dimexpr 0.0700\linewidth-2\tabcolsep\relax}>{\raggedright\arraybackslash}p{\dimexpr 0.8100\linewidth-2\tabcolsep\relax}>{\raggedright\arraybackslash}p{\dimexpr 0.1200\linewidth-2\tabcolsep\relax}}
\toprule
\textbf{\#} & \textbf{check} & \textbf{result} \\
\midrule\endfirsthead
\toprule
\textbf{\#} & \textbf{check} & \textbf{result} \\
\midrule\endhead
\endfoot
\bottomrule\endlastfoot
61 & \texttt{E\_\allowbreak{}H (16-component form, chiral orbitals) == gS/\allowbreak{}2 (κ + κ\textasciicircum{}*)\textasciicircum{}2 - gP/\allowbreak{}2 (κ - κ\textasciicircum{}*)\textasciicircum{}2 + (gV + gA)/\allowbreak{}2 (jL.jL + jR.jR) + (gV - gA) jL.jR} & True \\*
\end{longtable}}
\Needspace{6\baselineskip}\subsubsection*{9.3 The chiral Kohn–Sham equations}
{\footnotesize
\begin{longtable}{>{\raggedright\arraybackslash}p{\dimexpr 0.0700\linewidth-2\tabcolsep\relax}>{\raggedright\arraybackslash}p{\dimexpr 0.8100\linewidth-2\tabcolsep\relax}>{\raggedright\arraybackslash}p{\dimexpr 0.1200\linewidth-2\tabcolsep\relax}}
\toprule
\textbf{\#} & \textbf{check} & \textbf{result} \\
\midrule\endfirsthead
\toprule
\textbf{\#} & \textbf{check} & \textbf{result} \\
\midrule\endhead
\endfoot
\bottomrule\endlastfoot
62 & \texttt{chiral KS Hamiltonian: W\textasciicircum{}†.h\_\allowbreak{}KS(p).W == \{\{V\textasciicircum{}L\_\allowbreak{}0 - σW.(p + V\textasciicircum{}L), M\}, \{M\textasciicircum{}*, V\textasciicircum{}R\_\allowbreak{}0 + σW.(p + V\textasciicircum{}R)\}\}, M = m + S + I P, V\textasciicircum{}L = V - A, V\textasciicircum{}R = V + A} & True \\
63 & \texttt{Hartree (Wirtinger): M - m == ∂e\_\allowbreak{}H/\allowbreak{}∂κ == gS ρS + I gP ρP, V\textasciicircum{}L\_\allowbreak{}A == ∂e\_\allowbreak{}H/\allowbreak{}∂jL\textasciicircum{}A; V\textasciicircum{}L\_\allowbreak{}A == V\_\allowbreak{}A - A\_\allowbreak{}A and V\textasciicircum{}R\_\allowbreak{}A == V\_\allowbreak{}A + A\_\allowbreak{}A with V\_\allowbreak{}A = gV j\_\allowbreak{}A, A\_\allowbreak{}A = gA j5\_\allowbreak{}A} & True \\
64 & \texttt{chiral KS potentials are derivatives: ∂E\_\allowbreak{}H/\allowbreak{}∂φL\textasciicircum{}* == (V\textasciicircum{}L\_\allowbreak{}0 - σW[i] V\textasciicircum{}L\_\allowbreak{}i) φL + (M - m) φR and ∂E\_\allowbreak{}H/\allowbreak{}∂φR\textasciicircum{}* == (V\textasciicircum{}R\_\allowbreak{}0 + σW[i] V\textasciicircum{}R\_\allowbreak{}i) φR + (M - m)\textasciicircum{}* φL} & True \\*
\end{longtable}}
\Needspace{6\baselineskip}\subsubsection*{9.4 Chiral symmetry}
{\footnotesize
\begin{longtable}{>{\raggedright\arraybackslash}p{\dimexpr 0.0700\linewidth-2\tabcolsep\relax}>{\raggedright\arraybackslash}p{\dimexpr 0.8100\linewidth-2\tabcolsep\relax}>{\raggedright\arraybackslash}p{\dimexpr 0.1200\linewidth-2\tabcolsep\relax}}
\toprule
\textbf{\#} & \textbf{check} & \textbf{result} \\
\midrule\endfirsthead
\toprule
\textbf{\#} & \textbf{check} & \textbf{result} \\
\midrule\endhead
\endfoot
\bottomrule\endlastfoot
65 & \texttt{e\textasciicircum{}(Iθχ17) == cosθ ID16 - sinθ ω17 is REAL and orthogonal; U.EL == e\textasciicircum{}(-Iθ) EL, U.ER == e\textasciicircum{}(Iθ) ER} & True \\
66 & \texttt{U\textasciicircum{}†.βD.U == βD.(cos 2θ ID16 + I sin 2θ χ17); vector and axial densities invariant} & True \\
67 & \texttt{E\_\allowbreak{}H(e\textasciicircum{}(Iθχ17) ψ) - E\_\allowbreak{}H(ψ) == (gS - gP)/\allowbreak{}2 [(e\textasciicircum{}(4Iθ) - 1) κ\textasciicircum{}2 + (e\textasciicircum{}(-4Iθ) - 1) κ\textasciicircum{}*2]: U(1)\_\allowbreak{}A symmetric iff gP == gS} & True \\*
\end{longtable}}
\Needspace{6\baselineskip}\subsubsection*{9.5 Uniform matter in chiral form}
{\footnotesize
\begin{longtable}{>{\raggedright\arraybackslash}p{\dimexpr 0.0700\linewidth-2\tabcolsep\relax}>{\raggedright\arraybackslash}p{\dimexpr 0.8100\linewidth-2\tabcolsep\relax}>{\raggedright\arraybackslash}p{\dimexpr 0.1200\linewidth-2\tabcolsep\relax}}
\toprule
\textbf{\#} & \textbf{check} & \textbf{result} \\
\midrule\endfirsthead
\toprule
\textbf{\#} & \textbf{check} & \textbf{result} \\
\midrule\endhead
\endfoot
\bottomrule\endlastfoot
68 & \texttt{uniform matter in chiral form: κ == ρS/\allowbreak{}2, jL\textasciicircum{}0 == jR\textasciicircum{}0 == ρ/\allowbreak{}2, jL\textasciicircum{}i == jR\textasciicircum{}i == 0} & True \\
69 & \texttt{(h\_\allowbreak{}χ - V\_\allowbreak{}0)\textasciicircum{}2 == (p\textasciicircum{}2 + m*\textasciicircum{}2) ID16: eigenvalues V\_\allowbreak{}0 ± Sqrt[p\textasciicircum{}2 + m*\textasciicircum{}2], eight times each} & True \\*
\end{longtable}}
\Needspace{6\baselineskip}\subsubsection*{9.6 The chiral vacuum with the Dirac sea: spontaneous chiral symmetry breaking}
{\footnotesize
\begin{longtable}{>{\raggedright\arraybackslash}p{\dimexpr 0.0700\linewidth-2\tabcolsep\relax}>{\raggedright\arraybackslash}p{\dimexpr 0.8100\linewidth-2\tabcolsep\relax}>{\raggedright\arraybackslash}p{\dimexpr 0.1200\linewidth-2\tabcolsep\relax}}
\toprule
\textbf{\#} & \textbf{check} & \textbf{result} \\
\midrule\endfirsthead
\toprule
\textbf{\#} & \textbf{check} & \textbf{result} \\
\midrule\endhead
\endfoot
\bottomrule\endlastfoot
70 & \texttt{|S\textasciicircum{}6|/\allowbreak{}(2 Pi)\textasciicircum{}7 == 1/\allowbreak{}(120 Pi\textasciicircum{}4); J(0) == Λ\textasciicircum{}6/\allowbreak{}(720 Pi\textasciicircum{}4); g\_\allowbreak{}c == -1/\allowbreak{}(8 J(0)) == -90 Pi\textasciicircum{}4/\allowbreak{}Λ\textasciicircum{}6; J(M) -> J(0) and K(M) -> K(0) as M -> 0} & True \\
71 & \texttt{m = 0, g = 6/\allowbreak{}5 g\_\allowbreak{}c: the gap equation 1 + 8 g J(|M|) == 0 has the root |M| (32 digits); M == 2 g κ\textasciicircum{}* with κ == -4 M J; E is stationary there and lower than E(0)} & True \\
72 & \texttt{degenerate chiral vacua for m = 0: E(|M| e\textasciicircum{}(I ϕ)) == E(|M|) for ϕ = 1, 2, 4 (the phase is the Goldstone direction)} & True \\
73 & \texttt{explicit breaking m = 1/\allowbreak{}100: the real M == m - 8 g J(M) M solves the gap equation, E is stationary there, and every complex phase raises E} & True \\
74 & \texttt{the sea at one momentum: h\_\allowbreak{}χ\textasciicircum{}2 == E\textasciicircum{}2 ID16; Π\_\allowbreak{}- is a projector of rank 8; κ\_\allowbreak{}p == tr(Π\_\allowbreak{}- RL block) == -4 M\textasciicircum{}*/\allowbreak{}E; tr(Π\_\allowbreak{}-.h0) == -8E + 8 Re[(M - m) M\textasciicircum{}*]/\allowbreak{}E} & True \\*
\end{longtable}}
\Needspace{8\baselineskip}\subsection*{4.3 Independent Python checks: tools/ref\_dirac16.py}\addcontentsline{toc}{subsection}{4.3 Independent Python checks: tools/ref\_dirac16.py}
The program recomputes the constructions of sections 7–9 without Mathematica (Python 3.13.17, sympy 1.14.0, mpmath 1.3.0, numpy 2.5.3; its random test data come from the seed 71 of \texttt{random} and of numpy’s \texttt{default\_\allowbreak{}rng}). Its 90 boolean checks:

{\footnotesize
\begin{longtable}{>{\raggedright\arraybackslash}p{\dimexpr 0.0700\linewidth-2\tabcolsep\relax}>{\raggedright\arraybackslash}p{\dimexpr 0.8100\linewidth-2\tabcolsep\relax}>{\raggedright\arraybackslash}p{\dimexpr 0.1200\linewidth-2\tabcolsep\relax}}
\toprule
\textbf{\#} & \textbf{check} & \textbf{result} \\
\midrule\endfirsthead
\toprule
\textbf{\#} & \textbf{check} & \textbf{result} \\
\midrule\endhead
\endfoot
\bottomrule\endlastfoot
1 & \texttt{Gamma17 real, entries 0,+-1} & True \\
2 & \texttt{\{Gamma\_\allowbreak{}A,Gamma\_\allowbreak{}B\} = 2 eta17\_\allowbreak{}AB} & True \\
3 & \texttt{Gamma0 symmetric, Gamma\_\allowbreak{}i antisymmetric} & True \\
4 & \texttt{Gamma\textasciicircum{}A dagger = Gamma0 Gamma\textasciicircum{}A Gamma0} & True \\
5 & \texttt{alpha\textasciicircum{}i, beta real symmetric} & True \\
6 & \texttt{\{alpha\_\allowbreak{}i,alpha\_\allowbreak{}j\}=2delta, \{alpha\_\allowbreak{}i,beta\}=0, beta\textasciicircum{}2=1} & True \\
7 & \texttt{block forms alpha\textasciicircum{}i = [[0,tauO\_\allowbreak{}i],[-tauO\_\allowbreak{}i,0]], beta = [[0,1],[1,0]]} & True \\
8 & \texttt{omega = Gamma17[8] = [[0,-1],[1,0]], omega\textasciicircum{}2 = -1} & True \\
9 & \texttt{(i Dslash + m)(i Dslash - m) psi = -(box + m\textasciicircum{}2) psi} & True \\
10 & \texttt{Gamma0 (i d0 - h) psi = (i Dslash - m) psi} & True \\
11 & \texttt{h(p) = [[0,M],[M\textasciicircum{}T,0]]} & True \\
12 & \texttt{h(p)\textasciicircum{}2 = E\textasciicircum{}2 1} & True \\
13 & \texttt{M M\textasciicircum{}T = M\textasciicircum{}T M = E\textasciicircum{}2 1\_\allowbreak{}8} & True \\
14 & \texttt{h(p) u\_\allowbreak{}s = E u\_\allowbreak{}s,  h(-p) v\_\allowbreak{}s = -E v\_\allowbreak{}s} & True \\
15 & \texttt{(pslash - m) u\_\allowbreak{}s = 0, (pslash + m) v\_\allowbreak{}s = 0} & True \\
16 & \texttt{u\_\allowbreak{}s\textasciicircum{}T u\_\allowbreak{}r = 2E delta, v\_\allowbreak{}s\textasciicircum{}T v\_\allowbreak{}r = 2E delta} & True \\
17 & \texttt{u\_\allowbreak{}s(p)\textasciicircum{}T v\_\allowbreak{}r(p\textasciitilde{}) = 0} & True \\
18 & \texttt{sum\_\allowbreak{}s u\_\allowbreak{}s ubar\_\allowbreak{}s = pslash + m, sum\_\allowbreak{}s v\_\allowbreak{}s vbar\_\allowbreak{}s = pslash - m} & True \\
19 & \texttt{completeness sum u u\textasciicircum{}T + sum v(p\textasciitilde{}) v(p\textasciitilde{})\textasciicircum{}T = 2E 1\_\allowbreak{}16} & True \\
20 & \texttt{U\_\allowbreak{}p orthogonal and U\_\allowbreak{}p\textasciicircum{}T h(p) U\_\allowbreak{}p = diag(E 1\_\allowbreak{}8, -E 1\_\allowbreak{}8)} & True \\
21 & \texttt{charge conjugation: (pslash + m) omega u\_\allowbreak{}s = 0} & True \\
22 & \texttt{omega u\_\allowbreak{}s = sum\_\allowbreak{}r v\_\allowbreak{}r (-M\textasciicircum{}T/\allowbreak{}E)\_\allowbreak{}rs} & True \\
23 & \texttt{S Gamma\_\allowbreak{}K S\textasciicircum{}-1 = sum\_\allowbreak{}B Lambda\_\allowbreak{}BK Gamma\_\allowbreak{}B (boost)} & True \\
24 & \texttt{Lorentz covariance: (pslash(Lambda p) - m) S u\_\allowbreak{}s(p) = 0} & True \\
25 & \texttt{real matrices anticommuting with all Gamma17[A]: dimension 1 (rank 255)} & True \\
26 & \texttt{omega anticommutes with all Gamma17[A]} & True \\
27 & \texttt{chi = i omega: Hermitian, chi\textasciicircum{}2 = 1, \{chi,Gamma\_\allowbreak{}A\} = 0} & True \\
28 & \texttt{[chi, Sigma\_\allowbreak{}AB] = 0 (chirality is Lorentz invariant)} & True \\
29 & \texttt{P\_\allowbreak{}L, P\_\allowbreak{}R not real} & True \\
30 & \texttt{P\_\allowbreak{}L+P\_\allowbreak{}R = 1, P\_\allowbreak{}L\textasciicircum{}2 = P\_\allowbreak{}L, P\_\allowbreak{}R\textasciicircum{}2 = P\_\allowbreak{}R, P\_\allowbreak{}L P\_\allowbreak{}R = 0, rank 8 each} & True \\
31 & \texttt{Gamma\_\allowbreak{}A P\_\allowbreak{}L = P\_\allowbreak{}R Gamma\_\allowbreak{}A,  P\_\allowbreak{}L\textasciicircum{}dagger Gamma0 = Gamma0 P\_\allowbreak{}R} & True \\
32 & \texttt{E\_\allowbreak{}L, E\_\allowbreak{}R isometries: E\textasciicircum{}dag E = 1\_\allowbreak{}8, E\_\allowbreak{}L\textasciicircum{}dag E\_\allowbreak{}R = 0, E\_\allowbreak{}L E\_\allowbreak{}L\textasciicircum{}dag = P\_\allowbreak{}L, E\_\allowbreak{}R E\_\allowbreak{}R\textasciicircum{}dag = P\_\allowbreak{}R} & True \\
33 & \texttt{chi E\_\allowbreak{}L = -E\_\allowbreak{}L, chi E\_\allowbreak{}R = E\_\allowbreak{}R, W unitary} & True \\
34 & \texttt{W\textasciicircum{}dag h(p) W = [[-sigma.p, m],[m, sigma.p]]} & True \\
35 & \texttt{sigma\_\allowbreak{}i Hermitian, \{sigma\_\allowbreak{}i,sigma\_\allowbreak{}j\} = 2 delta} & True \\
36 & \texttt{chiral blocks: E\_\allowbreak{}L\textasciicircum{}dag beta E\_\allowbreak{}L = 0, E\_\allowbreak{}L\textasciicircum{}dag beta E\_\allowbreak{}R = 1, E\_\allowbreak{}L\textasciicircum{}dag alpha E\_\allowbreak{}L = -sigma, E\_\allowbreak{}R\textasciicircum{}dag alpha E\_\allowbreak{}R = sigma, E\_\allowbreak{}L\textasciicircum{}dag alpha E\_\allowbreak{}R = 0} & True \\
37 & \texttt{E\_\allowbreak{}L\textasciicircum{}dag Gamma0 (i Dslash - m) psi = (i d0 + tauO.d) phi\_\allowbreak{}L - m phi\_\allowbreak{}R} & True \\
38 & \texttt{E\_\allowbreak{}R\textasciicircum{}dag Gamma0 (i Dslash - m) psi = (i d0 - tauO.d) phi\_\allowbreak{}R - m phi\_\allowbreak{}L} & True \\
39 & \texttt{(i d0 - tauO.d)(i d0 + tauO.d) f = -box f (8 components)} & True \\
40 & \texttt{psi = E\_\allowbreak{}L phi\_\allowbreak{}L + E\_\allowbreak{}R phi\_\allowbreak{}R} & True \\
41 & \texttt{chiral form of the bilinears (beta i chi, chi, alpha chi)} & True \\
42 & \texttt{\{psi\_\allowbreak{}a(t,x), psi\_\allowbreak{}b\textasciicircum{}dag(t,y)\} = delta\_\allowbreak{}ab (1/\allowbreak{}V) sum\_\allowbreak{}q e\textasciicircum{}\{iq(x-y)\} (5 momenta, abstract mode algebra)} & True \\
43 & \texttt{\{psi\_\allowbreak{}a(t,x), psi\_\allowbreak{}b(t,y)\} = 0} & True \\
44 & \texttt{the three U\_\allowbreak{}k solve the Dirac equation (m = 1)} & True \\
45 & \texttt{Grassmann components anticommute: psi\_\allowbreak{}a psi\_\allowbreak{}b + psi\_\allowbreak{}b psi\_\allowbreak{}a = 0, psi\_\allowbreak{}a\textasciicircum{}2 = 0} & True \\
46 & \texttt{current conservation d\_\allowbreak{}A (psibar Gamma\textasciicircum{}A psi) = 0 (Grassmann field)} & True \\
47 & \texttt{dL\_\allowbreak{}int/\allowbreak{}dpsibar\_\allowbreak{}a = -[gS(psibar psi) psi + gP (psibar i chi psi) i chi psi + gV (psibar G\_\allowbreak{}A psi) G\textasciicircum{}A psi + gA (psibar G\_\allowbreak{}A chi psi) G\textasciicircum{}A chi psi]\_\allowbreak{}a} & True \\
48 & \texttt{h\_\allowbreak{}KS Hermitian for real S, P, V\_\allowbreak{}A, A\_\allowbreak{}A} & True \\
49 & \texttt{Gamma0 [pslash(eps,p) - m - S - i chi P - Gamma\textasciicircum{}A (V\_\allowbreak{}A + chi A\_\allowbreak{}A)] = eps - h\_\allowbreak{}KS(p)  (plane wave e\textasciicircum{}\{-i(eps t - p.x)\})} & True \\
50 & \texttt{W\textasciicircum{}dag h\_\allowbreak{}KS W = [[V\textasciicircum{}L\_\allowbreak{}0 - sigma.(p+V\textasciicircum{}L), M],[M\textasciicircum{}*, V\textasciicircum{}R\_\allowbreak{}0 + sigma.(p+V\textasciicircum{}R)]], V\textasciicircum{}L = V - A, V\textasciicircum{}R = V + A, M = m + S + iP} & True \\
51 & \texttt{Hartree potentials: dE\_\allowbreak{}H/\allowbreak{}dphi*\_\allowbreak{}a = (v\_\allowbreak{}KS phi)\_\allowbreak{}a with S = gS rhoS, P = gP rhoP, V\_\allowbreak{}A = gV j\_\allowbreak{}A, A\_\allowbreak{}A = gA j5\_\allowbreak{}A} & True \\
52 & \texttt{densities real: Gamma0 O Hermitian for O = 1, i chi, Gamma\textasciicircum{}A, Gamma\textasciicircum{}A chi} & True \\
53 & \texttt{chiral densities: rhoS = kappa + kappa\textasciicircum{}*, rhoP = i(kappa - kappa\textasciicircum{}*), j = jL + jR, j5 = jR - jL} & True \\
54 & \texttt{E\_\allowbreak{}H = gS/\allowbreak{}2 (k+k*)\textasciicircum{}2 - gP/\allowbreak{}2 (k-k*)\textasciicircum{}2 + (gV+gA)/\allowbreak{}2 (jL.jL + jR.jR) + (gV-gA) jL.jR} & True \\
55 & \texttt{V\textasciicircum{}L\_\allowbreak{}A = V\_\allowbreak{}A - A\_\allowbreak{}A, V\textasciicircum{}R\_\allowbreak{}A = V\_\allowbreak{}A + A\_\allowbreak{}A (Hartree, chiral currents)} & True \\
56 & \texttt{M - m = dE\_\allowbreak{}H/\allowbreak{}dkappa = gS(k+k*) - gP(k-k*) = S + iP} & True \\
57 & \texttt{dE\_\allowbreak{}H/\allowbreak{}dphi\_\allowbreak{}L* = (V\textasciicircum{}L\_\allowbreak{}0 - sigma\_\allowbreak{}i V\textasciicircum{}L\_\allowbreak{}i) phi\_\allowbreak{}L + (M - m) phi\_\allowbreak{}R ; dE\_\allowbreak{}H/\allowbreak{}dphi\_\allowbreak{}R* = (V\textasciicircum{}R\_\allowbreak{}0 + sigma\_\allowbreak{}i V\textasciicircum{}R\_\allowbreak{}i) phi\_\allowbreak{}R + (M - m)\textasciicircum{}* phi\_\allowbreak{}L} & True \\
58 & \texttt{U\_\allowbreak{}theta = cos(theta) - sin(theta) omega is real orthogonal; U E\_\allowbreak{}L = e\textasciicircum{}\{-i theta\} E\_\allowbreak{}L, U E\_\allowbreak{}R = e\textasciicircum{}\{i theta\} E\_\allowbreak{}R} & True \\
59 & \texttt{U\textasciicircum{}dag beta U = beta (cos 2theta + i sin 2theta chi);  vector and axial currents invariant} & True \\
60 & \texttt{E\_\allowbreak{}H(e\textasciicircum{}\{i theta chi\} psi) - E\_\allowbreak{}H(psi) = (gS-gP)/\allowbreak{}2 [(e\textasciicircum{}\{4i theta\}-1) k\textasciicircum{}2 + (e\textasciicircum{}\{-4i theta\}-1) k*\textasciicircum{}2]  ->  U(1)\_\allowbreak{}A symmetric iff gP = gS} & True \\
61 & \texttt{256 products Gamma\_\allowbreak{}K are orthogonal: Gamma\_\allowbreak{}K\textasciicircum{}-1 = Gamma\_\allowbreak{}K\textasciicircum{}T} & True \\
62 & \texttt{tr(Gamma\_\allowbreak{}K\textasciicircum{}-1 Gamma\_\allowbreak{}L) = 16 delta\_\allowbreak{}KL (256 x 256)} & True \\
63 & \texttt{Fierz completeness: sum\_\allowbreak{}K (Gamma\_\allowbreak{}K)\_\allowbreak{}ab (Gamma\_\allowbreak{}K\textasciicircum{}-1)\_\allowbreak{}cd = 16 delta\_\allowbreak{}ad delta\_\allowbreak{}cb} & True \\
64 & \texttt{delta E\_\allowbreak{}x = tr[Sigma\_\allowbreak{}x delta R] (random R, dR), Sigma\_\allowbreak{}x = -gS R - gP (i chi) R (i chi) - gV G\_\allowbreak{}A R G\textasciicircum{}A - gA G\_\allowbreak{}A chi R G\textasciicircum{}A chi} & True \\
65 & \texttt{E\_\allowbreak{}x = (1/\allowbreak{}2) tr[Sigma\_\allowbreak{}x R]} & True \\
66 & \texttt{E\_\allowbreak{}H = (gS/\allowbreak{}2)(tr R)\textasciicircum{}2 + ... = sum\_\allowbreak{}c (g\_\allowbreak{}c/\allowbreak{}2) tr(O\_\allowbreak{}c R) tr(O\textasciicircum{}c R) and dE\_\allowbreak{}H = tr[Sigma\_\allowbreak{}H dR]} & True \\
67 & \texttt{Fierz: tr[O1 R O2 R] = (1/\allowbreak{}16) sum\_\allowbreak{}K tr(Gamma\_\allowbreak{}K\textasciicircum{}-1 R) tr(O1 Gamma\_\allowbreak{}K O2 R) for all channel vertices} & True \\
68 & \texttt{uniform R: tr R = rho\_\allowbreak{}S, tr(Gamma\textasciicircum{}0 R) = rho} & True \\
69 & \texttt{(1/\allowbreak{}2) sum\_\allowbreak{}s [u\_\allowbreak{}s(p) ubar\_\allowbreak{}s(p) + u\_\allowbreak{}s(-p) ubar\_\allowbreak{}s(-p)]/\allowbreak{}(2E) = (E Gamma\_\allowbreak{}0 + m)/\allowbreak{}(2E)} & True \\
70 & \texttt{E\_\allowbreak{}x(uniform) = -gS/\allowbreak{}32 (rS\textasciicircum{}2+rho\textasciicircum{}2) + gP/\allowbreak{}32 (rS\textasciicircum{}2-rho\textasciicircum{}2) - gV/\allowbreak{}32 (8rS\textasciicircum{}2-6rho\textasciicircum{}2) + gA/\allowbreak{}32 (8rS\textasciicircum{}2+6rho\textasciicircum{}2)} & True \\
71 & \texttt{Sigma\_\allowbreak{}x(uniform) = S\_\allowbreak{}x + V\_\allowbreak{}x Gamma\_\allowbreak{}0, S\_\allowbreak{}x = rS/\allowbreak{}16(-gS+gP-8gV+8gA), V\_\allowbreak{}x = rho/\allowbreak{}16(-gS-gP+6gV+6gA)} & True \\
72 & \texttt{S\_\allowbreak{}x = dE\_\allowbreak{}x/\allowbreak{}drho\_\allowbreak{}S, V\_\allowbreak{}x = dE\_\allowbreak{}x/\allowbreak{}drho} & True \\
73 & \texttt{uniform Hartree-Fock = Hartree with G\_\allowbreak{}S = gS + (-gS+gP-8gV+8gA)/\allowbreak{}16, G\_\allowbreak{}V = gV + (-gS-gP+6gV+6gA)/\allowbreak{}16} & True \\
74 & \texttt{uniform matter: Hartree terms of P and A vanish (rho\_\allowbreak{}P = 0, j5 = 0, j\textasciicircum{}i = 0)} & True \\
75 & \texttt{Wick (3 particles, 16 modes): <:(psibar O psi)(psibar O' psi):> = tr(O R) tr(O' R) - tr(O R O' R) for all 18 vertex pairs} & True \\
76 & \texttt{c7 = 8 |S\textasciicircum{}6|/\allowbreak{}(2pi)\textasciicircum{}7 = 1/\allowbreak{}(15 pi\textasciicircum{}4) and rho = c7 kF\textasciicircum{}7/\allowbreak{}7 = kF\textasciicircum{}7/\allowbreak{}(105 pi\textasciicircum{}4)} & True \\
77 & \texttt{Hartree gap: m* - m - gS rhoS(m*) = 0} & True \\
78 & \texttt{m* is a stationary point of eps(m*) = ekin(m*) - (m*-m)\textasciicircum{}2/\allowbreak{}(2gS) + gV rho\textasciicircum{}2/\allowbreak{}2} & True \\
79 & \texttt{total energy = sum eps\_\allowbreak{}k - E\_\allowbreak{}H (double counting) = ekin - (m*-m) rhoS + E\_\allowbreak{}H} & True \\
80 & \texttt{Hugenholtz-Van Hove at saturation: E/\allowbreak{}A = eps\_\allowbreak{}F = d eps/\allowbreak{}d rho} & True \\
81 & \texttt{d(eps/\allowbreak{}rho)/\allowbreak{}dkF = 0 at kF0 (minimum: E/\allowbreak{}A(0.95 kF0), E/\allowbreak{}A(1.05 kF0) larger)} & True \\
82 & \texttt{Hugenholtz-Van Hove mu = d eps/\allowbreak{}d rho = eps\_\allowbreak{}F at kF = 1 (not saturated)} & True \\
83 & \texttt{HF self-consistency with the exchange self-energy: m* - m = gS rhoS + S\_\allowbreak{}x, V0 = gV rho + V\_\allowbreak{}x} & True \\
84 & \texttt{HF energy: ekin - (m*-m) rhoS + E\_\allowbreak{}H + E\_\allowbreak{}x = eps(G\_\allowbreak{}S, G\_\allowbreak{}V)} & True \\
85 & \texttt{HF Hugenholtz-Van Hove: d eps/\allowbreak{}d rho = eps\_\allowbreak{}F at kF = 3/\allowbreak{}4} & True \\
86 & \texttt{J(0) = Lambda\textasciicircum{}6/\allowbreak{}(720 pi\textasciicircum{}4), g\_\allowbreak{}c = -1/\allowbreak{}(8 J(0)) = -90 pi\textasciicircum{}4/\allowbreak{}Lambda\textasciicircum{}6} & True \\
87 & \texttt{NJL m = 0: gap 1 + 8 g J(|M|) = 0 solved, M = 2 g kappa\textasciicircum{}*, stationary and lower than M = 0} & True \\
88 & \texttt{NJL m = 0: E(|M| e\textasciicircum{}\{i theta\}) independent of theta (degenerate chiral vacua)} & True \\
89 & \texttt{NJL m = 1/\allowbreak{}100: M = m - 8 g J(M) M real, stationary; E(M e\textasciicircum{}\{i theta\}) > E(M) for theta != 0} & True \\
90 & \texttt{sea at one momentum: eigenvalues -E (x8), +E (x8); kappa\_\allowbreak{}p = tr[(Pi\_\allowbreak{}-)\_\allowbreak{}\{RL\}] = -4 M\textasciicircum{}*/\allowbreak{}E; tr[Pi\_\allowbreak{}- h\_\allowbreak{}0] = -8E + 8 Re[(M-m) M\textasciicircum{}*]/\allowbreak{}E} & True \\*
\end{longtable}}
The numbers it computes (30-digit arithmetic). \texttt{tools/\allowbreak{}verify\_\allowbreak{}7p1\_\allowbreak{}nb.py}, whose output is printed below together with that of \texttt{tools/\allowbreak{}ref\_\allowbreak{}dirac16.py}, compares them with the numbers printed by 7+1.nb:

{\footnotesize
\begin{longtable}{>{\raggedright\arraybackslash}p{\dimexpr 0.5789\linewidth-2\tabcolsep\relax}>{\raggedright\arraybackslash}p{\dimexpr 0.4211\linewidth-2\tabcolsep\relax}}
\toprule
\textbf{quantity} & \textbf{value} \\
\midrule\endfirsthead
\toprule
\textbf{quantity} & \textbf{value} \\
\midrule\endhead
\endfoot
\bottomrule\endlastfoot
\texttt{Hartree saturation kF0} & $0.75663615070149289469$ \\
\texttt{Hartree rho0} & $0.000013881002132316877197$ \\
\texttt{Hartree m*} & $0.34809734014624326514$ \\
\texttt{Hartree rhoS} & $6.5190265985375673486e-6$ \\
\texttt{Hartree V0} & $0.13881002132316877197$ \\
\texttt{Hartree E/\allowbreak{}A - m} & $-0.028321427964738575854$ \\
\texttt{Hartree eps\_\allowbreak{}F} & $0.97167857203526142415$ \\
\texttt{Hartree kF=1: m*} & $0.080964430277606126329$ \\
\texttt{Hartree kF=1: E/\allowbreak{}A - m} & $0.41086484863433217115$ \\
\texttt{Hartree kF=1: eps\_\allowbreak{}F} & $1.9809848613081542876$ \\
\texttt{HF couplings G\_\allowbreak{}S, G\_\allowbreak{}V} & $(mpf('-98750.0'), mpf('20000.0'))$ \\
\texttt{HF kF=3/\allowbreak{}4: m*} & $0.36643630198154392124$ \\
\texttt{HF kF=3/\allowbreak{}4: V0} & $0.26101775473161552570$ \\
\texttt{HF kF=3/\allowbreak{}4: E/\allowbreak{}A - m} & $0.039146338016119861526$ \\
\texttt{Hartree kF=3/\allowbreak{}4: m*} & $0.36297394764803019822$ \\
\texttt{Hartree kF=3/\allowbreak{}4: E/\allowbreak{}A - m} & $-0.028065375846872960879$ \\
\texttt{NJL g\_\allowbreak{}c} & $-8766.8181930602193513$ \\
\texttt{NJL g = 6/\allowbreak{}5 g\_\allowbreak{}c} & $-10520.181831672263222$ \\
\texttt{NJL |M| (m = 0)} & $0.55369385182924169122$ \\
\texttt{NJL kappa = -4 M J(|M|) (M real)} & $-0.000026315792858364874200$ \\
\texttt{NJL E(|M|) - E(0)} & $-1.3102449629010022473e-6$ \\
\texttt{NJL M (m = 1/\allowbreak{}100)} & $0.58531284674622051123$ \\*
\end{longtable}}
\textbf{tools/verify\_7p1\_nb.py (parses the evaluated 7+1.nb; exit code 0)}

\begin{Verbatim}[fontsize=\small,frame=single,framesep=2mm]
input cells                                                                                    
    ...: 189
output cells                                                                                   
    ...: 554
message cells (must be 0)                                                                      
    ...: 0
check outputs                                                                                  
    ...: 443
check outputs that are True                                                                    
    ...: 443
check["..."] calls in the input cells                                                          
    ...: 443
$Version output                                                                                
    ...: ['{"15.0.1 for Microsoft Windows (64-bit) (July 2, 2026)","Notebook$$75$191583`"}']
private notebook context                                                                       
    ...: True
Length[checks]                                                                                 
    ...: 479
failedChecks == {}                                                                             
    ...: True
allChecksPassed output True                                                                    
    ...: True
tau[A] printed == 4+4 tau[A] (python)                                                          
    ...: True
tauO[A] printed == python tauO[A]                                                              
    ...: True
Gamma17[A] printed == python G17[A]                                                            
    ...: True
gridO printed == python octonion table                                                         
    ...: True
gridx44 printed == python split table (4+4.nb)                                                 
    ...: True
Lambda boost/rotation printed                                                                  
    ...: True
8 real on-shell spinors printed solve (pslash - m) w = 0 (python check)                        
    ...: True
classification table printed                                                                   
    ...: True
section 7: check outputs (all True)                                                            
    ...: 35 (True)
section 8: check outputs (all True)                                                            
    ...: 23 (True)
section 9: check outputs (all True)                                                            
    ...: 16 (True)
section 7: printed Sqrt[5] u_s(p) == python (5 e_s; M^T e_s) and (pslash - m) u_s == 0         
    ...: True
section 8: printed closed forms of int k^6 m*/E and int k^6 E: d/dkF == integrands, zero at kF
    ...= 0 (sympy)                                                                             
    ...: True
sections 8-9: 23 printed numbers agree with tools/ref_dirac16.py to 17 digits                  
    ...: True
sections 8-9: G_S == -98750, G_V == 20000 printed                                              
    ...: True
tools/ref_dirac16.py: all its boolean checks True                                              
    ...: True
sections 0-6: the 166 input and text cells of the backup are all present, 159 unchanged, 7
    ...changed and 1 added by the signature explanation of 27 October 2026                     
    ...: True
section 11: check outputs (all True)                                                           
    ...: 7 (True)
section 12: check outputs (all True)                                                           
    ...: 28 (True)
section 13: check outputs (all True)                                                           
    ...: 15 (True)
section 14: check outputs (all True)                                                           
    ...: 14 (True)
section 15: check outputs (all True)                                                           
    ...: 16 (True)
section 16: check outputs (all True)                                                           
    ...: 8 (True)
section 17: check outputs (all True)                                                           
    ...: 11 (True)
section 18: check outputs (all True)                                                           
    ...: 7 (True)
section 19: check outputs (all True)                                                           
    ...: 6 (True)
summary table: entries == Length[checks], all True                                             
    ...: True
summary table: every printed check is listed                                                   
    ...: True
sections 0-10: the 368 input and text cells of the backup taken before sections 11-20 are all
    ...present, 360 unchanged, 8 changed and 1 added by the signature explanation of 27 October
    ...2026                                                                                    
    ...: True
sections 0-20: the 601 input and text cells of the backup taken before the signature
    ...explanation are all present, 586 unchanged, 15 changed and 4 added by that explanation  
    ...: True
sections 0-20, outputs: of the 382 outputs of the backup taken before the signature explanation
    ...373 are identical; the others differ only in the check counts and tables of the summary
    ...sections, the declared list {7,1,0}, the private context name and measured times; 5 new
    ...check outputs, all True : True
sections 14-19: 27 groups of printed numbers agree with the Python programs (relative 1e-6;
    ...rounded prints 1e-4; ODE and ill-conditioned results 1e-5)                              
    ...: True
section 20: check outputs (all True)                                                           
    ...: 14 (True)
section 21: check outputs (all True)                                                           
    ...: 19 (True)
section 22: check outputs (all True)                                                           
    ...: 9 (True)
section 23: check outputs (all True)                                                           
    ...: 6 (True)
stage 4: without sections 20-23 (and with the old title of the summary section) the 605 input
    ...and text cells equal those of the backup taken before sections 20-23 were appended      
    ...: True
stage 4: the 387 outputs of sections 0-19 and of the summary equal those of the backup (summary
    ...counts and tables, the private context name and measured times excepted)                
    ...: True
sections 20-23: 34 groups of printed numbers agree with tools/ref_coupled.py (tolerances per
    ...group, from 1e-14 for exact constants to 2e-3 for numbers read off numerical solutions) 
    ...: True
section 24: check outputs (all True)                                                           
    ...: 20 (True)
stage 5: without section 24 (and with the old title of the summary section) the 732 input and
    ...text cells equal those of the backup taken before section 24 was appended, after undoing
    ...the 1 documented correction(s) of section 23.3                                          
    ...: True
stage 5: the 453 outputs of sections 0-23 and of the summary equal those of the backup (summary
    ...counts and tables, the private context name and measured times excepted; 1 output(s)
    ...differ only by the documented correction)                                               
    ...: True
section 24 numbers: 29 groups of printed numbers agree with tools/ref_dark.py (tolerances per
    ...group, from 1e-14 for closed forms to 1e-5 for a number read off a numerical solution)  
    ...: True
section 25: check outputs (all True)                                                           
    ...: 27 (True)
stage 6: without section 25 (and with the old title of the summary section) the 786 input and
    ...text cells equal those of the backup taken before section 25 was appended               
    ...: True
stage 6: the 481 outputs of sections 0-24 and of the summary equal those of the backup (summary
    ...counts and tables, the private context name and measured times excepted)                
    ...: True
section 25 numbers: 17 groups of printed numbers agree with tools/ref_torus_brane.py (integers
    ...and exact values exactly, floating-point numbers to 1e-15 ... 1e-6 per group)           
    ...: True
section 26: check outputs (all True)                                                           
    ...: 33 (True)
stage 7: without section 26 (and with the old title of the summary section) the 815 input and
    ...text cells equal those of the backup taken before section 26 was appended               
    ...: True
stage 7: the 515 outputs of sections 0-25 and of the summary equal those of the backup (summary
    ...counts and tables, the private context name and measured times excepted)                
    ...: True
section 26 numbers: 15 groups of printed numbers agree with tools/ref_radion.py (integers and
    ...exact values exactly, floating-point numbers to 1e-15 ... 1e-9 per group)               
    ...: True

VERDICT: ALL CROSS-CHECKS PASSED
\end{Verbatim}
\textbf{tools/ref\_dirac16.py (independent checks of sections 7-9; exit code 0)}

\begin{Verbatim}[fontsize=\small,frame=single,framesep=2mm]
Gamma17 real, entries 0,+-1 : True
{Gamma_A,Gamma_B} = 2 eta17_AB : True
Gamma0 symmetric, Gamma_i antisymmetric : True
Gamma^A dagger = Gamma0 Gamma^A Gamma0 : True
alpha^i, beta real symmetric : True
{alpha_i,alpha_j}=2delta, {alpha_i,beta}=0, beta^2=1 : True
block forms alpha^i = [[0,tauO_i],[-tauO_i,0]], beta = [[0,1],[1,0]] : True
omega = Gamma17[8] = [[0,-1],[1,0]], omega^2 = -1 : True
(i Dslash + m)(i Dslash - m) psi = -(box + m^2) psi : True
Gamma0 (i d0 - h) psi = (i Dslash - m) psi : True
h(p) = [[0,M],[M^T,0]] : True
h(p)^2 = E^2 1 : True
M M^T = M^T M = E^2 1_8 : True
h(p) u_s = E u_s,  h(-p) v_s = -E v_s : True
(pslash - m) u_s = 0, (pslash + m) v_s = 0 : True
u_s^T u_r = 2E delta, v_s^T v_r = 2E delta : True
u_s(p)^T v_r(p~) = 0 : True
sum_s u_s ubar_s = pslash + m, sum_s v_s vbar_s = pslash - m : True
completeness sum u u^T + sum v(p~) v(p~)^T = 2E 1_16 : True
U_p orthogonal and U_p^T h(p) U_p = diag(E 1_8, -E 1_8) : True
charge conjugation: (pslash + m) omega u_s = 0 : True
omega u_s = sum_r v_r (-M^T/E)_rs : True
S Gamma_K S^-1 = sum_B Lambda_BK Gamma_B (boost) : True
Lorentz covariance: (pslash(Lambda p) - m) S u_s(p) = 0 : True
real matrices anticommuting with all Gamma17[A]: dimension 1 (rank 255) : True
omega anticommutes with all Gamma17[A] : True
chi = i omega: Hermitian, chi^2 = 1, {chi,Gamma_A} = 0 : True
[chi, Sigma_AB] = 0 (chirality is Lorentz invariant) : True
P_L, P_R not real : True
P_L+P_R = 1, P_L^2 = P_L, P_R^2 = P_R, P_L P_R = 0, rank 8 each : True
Gamma_A P_L = P_R Gamma_A,  P_L^dagger Gamma0 = Gamma0 P_R : True
E_L, E_R isometries: E^dag E = 1_8, E_L^dag E_R = 0, E_L E_L^dag = P_L, E_R E_R^dag = P_R : True
chi E_L = -E_L, chi E_R = E_R, W unitary : True
W^dag h(p) W = [[-sigma.p, m],[m, sigma.p]] : True
sigma_i Hermitian, {sigma_i,sigma_j} = 2 delta : True
chiral blocks: E_L^dag beta E_L = 0, E_L^dag beta E_R = 1, E_L^dag alpha E_L = -sigma, E_R^dag
    ...alpha E_R = sigma, E_L^dag alpha E_R = 0 : True
E_L^dag Gamma0 (i Dslash - m) psi = (i d0 + tauO.d) phi_L - m phi_R : True
E_R^dag Gamma0 (i Dslash - m) psi = (i d0 - tauO.d) phi_R - m phi_L : True
(i d0 - tauO.d)(i d0 + tauO.d) f = -box f (8 components) : True
psi = E_L phi_L + E_R phi_R : True
chiral form of the bilinears (beta i chi, chi, alpha chi) : True
{psi_a(t,x), psi_b^dag(t,y)} = delta_ab (1/V) sum_q e^{iq(x-y)} (5 momenta, abstract mode
    ...algebra) : True
{psi_a(t,x), psi_b(t,y)} = 0 : True
the three U_k solve the Dirac equation (m = 1) : True
Grassmann components anticommute: psi_a psi_b + psi_b psi_a = 0, psi_a^2 = 0 : True
current conservation d_A (psibar Gamma^A psi) = 0 (Grassmann field) : True
dL_int/dpsibar_a = -[gS(psibar psi) psi + gP (psibar i chi psi) i chi psi + gV (psibar G_A psi)
    ...G^A psi + gA (psibar G_A chi psi) G^A chi psi]_a : True
h_KS Hermitian for real S, P, V_A, A_A : True
Gamma0 [pslash(eps,p) - m - S - i chi P - Gamma^A (V_A + chi A_A)] = eps - h_KS(p)  (plane wave
    ...e^{-i(eps t - p.x)}) : True
W^dag h_KS W = [[V^L_0 - sigma.(p+V^L), M],[M^*, V^R_0 + sigma.(p+V^R)]], V^L = V - A, V^R = V
    ...+ A, M = m + S + iP : True
Hartree potentials: dE_H/dphi*_a = (v_KS phi)_a with S = gS rhoS, P = gP rhoP, V_A = gV j_A,
    ...A_A = gA j5_A : True
densities real: Gamma0 O Hermitian for O = 1, i chi, Gamma^A, Gamma^A chi : True
chiral densities: rhoS = kappa + kappa^*, rhoP = i(kappa - kappa^*), j = jL + jR, j5 = jR - jL
    ...: True
E_H = gS/2 (k+k*)^2 - gP/2 (k-k*)^2 + (gV+gA)/2 (jL.jL + jR.jR) + (gV-gA) jL.jR : True
V^L_A = V_A - A_A, V^R_A = V_A + A_A (Hartree, chiral currents) : True
M - m = dE_H/dkappa = gS(k+k*) - gP(k-k*) = S + iP : True
dE_H/dphi_L* = (V^L_0 - sigma_i V^L_i) phi_L + (M - m) phi_R ; dE_H/dphi_R* = (V^R_0 + sigma_i
    ...V^R_i) phi_R + (M - m)^* phi_L : True
U_theta = cos(theta) - sin(theta) omega is real orthogonal; U E_L = e^{-i theta} E_L, U E_R =
    ...e^{i theta} E_R : True
U^dag beta U = beta (cos 2theta + i sin 2theta chi);  vector and axial currents invariant : True
E_H(e^{i theta chi} psi) - E_H(psi) = (gS-gP)/2 [(e^{4i theta}-1) k^2 + (e^{-4i theta}-1) k*^2]
    ...->  U(1)_A symmetric iff gP = gS : True
256 products Gamma_K are orthogonal: Gamma_K^-1 = Gamma_K^T : True
tr(Gamma_K^-1 Gamma_L) = 16 delta_KL (256 x 256) : True
Fierz completeness: sum_K (Gamma_K)_ab (Gamma_K^-1)_cd = 16 delta_ad delta_cb : True
delta E_x = tr[Sigma_x delta R] (random R, dR), Sigma_x = -gS R - gP (i chi) R (i chi) - gV G_A
    ...R G^A - gA G_A chi R G^A chi : True
E_x = (1/2) tr[Sigma_x R] : True
E_H = (gS/2)(tr R)^2 + ... = sum_c (g_c/2) tr(O_c R) tr(O^c R) and dE_H = tr[Sigma_H dR] : True
Fierz: tr[O1 R O2 R] = (1/16) sum_K tr(Gamma_K^-1 R) tr(O1 Gamma_K O2 R) for all channel
    ...vertices : True
uniform R: tr R = rho_S, tr(Gamma^0 R) = rho : True
(1/2) sum_s [u_s(p) ubar_s(p) + u_s(-p) ubar_s(-p)]/(2E) = (E Gamma_0 + m)/(2E) : True
E_x(uniform) = -gS/32 (rS^2+rho^2) + gP/32 (rS^2-rho^2) - gV/32 (8rS^2-6rho^2) + gA/32
    ...(8rS^2+6rho^2) : True
Sigma_x(uniform) = S_x + V_x Gamma_0, S_x = rS/16(-gS+gP-8gV+8gA), V_x = rho/16(-gS-gP+6gV+6gA)
    ...: True
S_x = dE_x/drho_S, V_x = dE_x/drho : True
uniform Hartree-Fock = Hartree with G_S = gS + (-gS+gP-8gV+8gA)/16, G_V = gV +
    ...(-gS-gP+6gV+6gA)/16 : True
uniform matter: Hartree terms of P and A vanish (rho_P = 0, j5 = 0, j^i = 0) : True
Wick (3 particles, 16 modes): <:(psibar O psi)(psibar O' psi):> = tr(O R) tr(O' R) - tr(O R O'
    ...R) for all 18 vertex pairs : True
c7 = 8 |S^6|/(2pi)^7 = 1/(15 pi^4) and rho = c7 kF^7/7 = kF^7/(105 pi^4) : True
NUM Hartree saturation kF0 : 0.75663615070149289469
NUM Hartree rho0 : 0.000013881002132316877197
NUM Hartree m* : 0.34809734014624326514
NUM Hartree rhoS : 6.5190265985375673486e-6
NUM Hartree V0 : 0.13881002132316877197
NUM Hartree E/A - m : -0.028321427964738575854
NUM Hartree eps_F : 0.97167857203526142415
Hartree gap: m* - m - gS rhoS(m*) = 0 : True
m* is a stationary point of eps(m*) = ekin(m*) - (m*-m)^2/(2gS) + gV rho^2/2 : True
total energy = sum eps_k - E_H (double counting) = ekin - (m*-m) rhoS + E_H : True
Hugenholtz-Van Hove at saturation: E/A = eps_F = d eps/d rho : True
d(eps/rho)/dkF = 0 at kF0 (minimum: E/A(0.95 kF0), E/A(1.05 kF0) larger) : True
Hugenholtz-Van Hove mu = d eps/d rho = eps_F at kF = 1 (not saturated) : True
NUM Hartree kF=1: m* : 0.080964430277606126329
NUM Hartree kF=1: E/A - m : 0.41086484863433217115
NUM Hartree kF=1: eps_F : 1.9809848613081542876
NUM HF couplings G_S, G_V : -98750.0, 20000.0
HF self-consistency with the exchange self-energy: m* - m = gS rhoS + S_x, V0 = gV rho + V_x :
    ...True
HF energy: ekin - (m*-m) rhoS + E_H + E_x = eps(G_S, G_V) : True
HF Hugenholtz-Van Hove: d eps/d rho = eps_F at kF = 3/4 : True
NUM HF kF=3/4: m* : 0.36643630198154392124
NUM HF kF=3/4: V0 : 0.2610177547316155257
NUM HF kF=3/4: E/A - m : 0.039146338016119861526
NUM Hartree kF=3/4: m* : 0.36297394764803019822
NUM Hartree kF=3/4: E/A - m : -0.028065375846872960879
J(0) = Lambda^6/(720 pi^4), g_c = -1/(8 J(0)) = -90 pi^4/Lambda^6 : True
NUM NJL g_c : -8766.8181930602193513
NUM NJL g = 6/5 g_c : -10520.181831672263222
NUM NJL |M| (m = 0) : 0.55369385182924169122
NUM NJL kappa = -4 M J(|M|) (M real) : -0.0000263157928583648742
NUM NJL E(|M|) - E(0) : -1.3102449629010022473e-6
NJL m = 0: gap 1 + 8 g J(|M|) = 0 solved, M = 2 g kappa^*, stationary and lower than M = 0 :
    ...True
NJL m = 0: E(|M| e^{i theta}) independent of theta (degenerate chiral vacua) : True
NUM NJL M (m = 1/100) : 0.58531284674622051123
NJL m = 1/100: M = m - 8 g J(M) M real, stationary; E(M e^{i theta}) > E(M) for theta != 0 :
    ...True
sea at one momentum: eigenvalues -E (x8), +E (x8); kappa_p = tr[(Pi_-)_{RL}] = -4 M^*/E;
    ...tr[Pi_- h_0] = -8E + 8 Re[(M-m) M^*]/E : True
boolean checks: 90   failed: 0
VERDICT: ALL CHECKS PASSED
\end{Verbatim}
\Needspace{8\baselineskip}\subsection*{4.4 Hashes of the files}\addcontentsline{toc}{subsection}{4.4 Hashes of the files}
{\footnotesize
\begin{longtable}{>{\raggedright\arraybackslash}p{\dimexpr 0.3878\linewidth-2\tabcolsep\relax}>{\raggedright\arraybackslash}p{\dimexpr 0.6122\linewidth-2\tabcolsep\relax}}
\toprule
\textbf{file} & \textbf{SHA-256} \\
\midrule\endfirsthead
\toprule
\textbf{file} & \textbf{SHA-256} \\
\midrule\endhead
\endfoot
\bottomrule\endlastfoot
\texttt{7+1.nb} (evaluated) & \texttt{41d252f9\allowbreak{}852f210c\allowbreak{}228f45e1\allowbreak{}ac956d2d\allowbreak{}7ef26fae\allowbreak{}b0062a58\allowbreak{}80280e1a\allowbreak{}ae0effc7} \\
\texttt{4+4.nb} & \texttt{74002323\allowbreak{}3eaa0cd2\allowbreak{}4ca0e2e7\allowbreak{}01009467\allowbreak{}1b1d314b\allowbreak{}698114c0\allowbreak{}8db94050\allowbreak{}d3e0aed5} \\
\texttt{7+1\_\allowbreak{}2026-10-26\_\allowbreak{}before\_\allowbreak{}dirac16complex.nb} (in backups/, workspace only) & \texttt{f7abe727\allowbreak{}9b5540e9\allowbreak{}14de876d\allowbreak{}c055c199\allowbreak{}bd491b9c\allowbreak{}2d9522bf\allowbreak{}28dbc236\allowbreak{}89e237ce} \\
\texttt{tools/\allowbreak{}cayley\_\allowbreak{}dickson\_\allowbreak{}formulas.py} & \texttt{bdde8c4c\allowbreak{}8199b09e\allowbreak{}f36ae4e1\allowbreak{}5d5d53bc\allowbreak{}62744fc8\allowbreak{}5b4c5ba5\allowbreak{}2eab5516\allowbreak{}f77b5d5f} \\
\texttt{tools/\allowbreak{}cayley\_\allowbreak{}dickson\_\allowbreak{}match.py} & \texttt{73f2cc58\allowbreak{}4bbe1a36\allowbreak{}20c50890\allowbreak{}ffe5fb29\allowbreak{}9857b04a\allowbreak{}14bc34fe\allowbreak{}d2f9eaef\allowbreak{}d2e36c4e} \\
\texttt{tools/\allowbreak{}classify\_\allowbreak{}clifford.py} & \texttt{42291276\allowbreak{}f46b0461\allowbreak{}83fdbb7b\allowbreak{}62be7c5c\allowbreak{}edfd695f\allowbreak{}e9da96ac\allowbreak{}256afa23\allowbreak{}580930a9} \\
\texttt{tools/\allowbreak{}doc\_\allowbreak{}brane.py} & \texttt{a39b4bf8\allowbreak{}0b5ad756\allowbreak{}194060ec\allowbreak{}e3fdd94e\allowbreak{}b604bf14\allowbreak{}e441984c\allowbreak{}64824736\allowbreak{}ca066b99} \\
\texttt{tools/\allowbreak{}doc\_\allowbreak{}brane\_\allowbreak{}provenance.py} & \texttt{41d2de4c\allowbreak{}4d3c9553\allowbreak{}62ad7f48\allowbreak{}3813785b\allowbreak{}7c686a24\allowbreak{}782ae34a\allowbreak{}135983c4\allowbreak{}05d3afe6} \\
\texttt{tools/\allowbreak{}doc\_\allowbreak{}brane\_\allowbreak{}symbols.py} & \texttt{31ecfae0\allowbreak{}0a2dc049\allowbreak{}33b0e057\allowbreak{}59057f4d\allowbreak{}caccee64\allowbreak{}238b3272\allowbreak{}eeedb930\allowbreak{}1edb85ef} \\
\texttt{tools/\allowbreak{}doc\_\allowbreak{}common.py} & \texttt{51995b29\allowbreak{}9ba2aa09\allowbreak{}47fcbfd5\allowbreak{}93486cb9\allowbreak{}f2cd724d\allowbreak{}15f7a430\allowbreak{}35f7a47f\allowbreak{}8ade454c} \\
\texttt{tools/\allowbreak{}doc\_\allowbreak{}coupled.py} & \texttt{1c43ad45\allowbreak{}64093f14\allowbreak{}040ce792\allowbreak{}b69635ef\allowbreak{}287c055b\allowbreak{}26ae540c\allowbreak{}bba7e40d\allowbreak{}806c49dd} \\
\texttt{tools/\allowbreak{}doc\_\allowbreak{}coupled\_\allowbreak{}provenance.py} & \texttt{b41808d1\allowbreak{}447e1242\allowbreak{}2940203d\allowbreak{}9301ae42\allowbreak{}32482691\allowbreak{}5bbf9697\allowbreak{}4b8d6bb9\allowbreak{}793fa972} \\
\texttt{tools/\allowbreak{}doc\_\allowbreak{}coupled\_\allowbreak{}symbols.py} & \texttt{e0557a52\allowbreak{}49c59e43\allowbreak{}087326d2\allowbreak{}d5d9f779\allowbreak{}37a38bf5\allowbreak{}2845374d\allowbreak{}aca4e004\allowbreak{}85981f82} \\
\texttt{tools/\allowbreak{}doc\_\allowbreak{}dark.py} & \texttt{cd2606f5\allowbreak{}caacad55\allowbreak{}ae77213a\allowbreak{}bbe33ab2\allowbreak{}fc89ee33\allowbreak{}c2049a8d\allowbreak{}44b05b5b\allowbreak{}2606bf9e} \\
\texttt{tools/\allowbreak{}doc\_\allowbreak{}dark\_\allowbreak{}provenance.py} & \texttt{89956664\allowbreak{}f5d6992c\allowbreak{}054917d4\allowbreak{}52c3e973\allowbreak{}ae8c1f36\allowbreak{}001fe83a\allowbreak{}5fb7e8c6\allowbreak{}f8f2ea14} \\
\texttt{tools/\allowbreak{}doc\_\allowbreak{}dark\_\allowbreak{}symbols.py} & \texttt{ee662066\allowbreak{}ae88d927\allowbreak{}c4d85991\allowbreak{}256e845b\allowbreak{}67b3e08e\allowbreak{}a5d6ba94\allowbreak{}06bb1d42\allowbreak{}1c0d3481} \\
\texttt{tools/\allowbreak{}doc\_\allowbreak{}dirac16.py} & \texttt{62811d8c\allowbreak{}383ce299\allowbreak{}894858cc\allowbreak{}57fe3cd7\allowbreak{}80088921\allowbreak{}3f77e83d\allowbreak{}1aa161f6\allowbreak{}af206070} \\
\texttt{tools/\allowbreak{}doc\_\allowbreak{}dirac16\_\allowbreak{}common.py} & \texttt{9324b4fa\allowbreak{}e59e5101\allowbreak{}9125ebe0\allowbreak{}9157f3c8\allowbreak{}c9489560\allowbreak{}e1383a69\allowbreak{}7fe6c22a\allowbreak{}3a93e208} \\
\texttt{tools/\allowbreak{}doc\_\allowbreak{}dirac16\_\allowbreak{}provenance.py} & \texttt{762f8594\allowbreak{}ea24910b\allowbreak{}fb402e40\allowbreak{}8b80e589\allowbreak{}014c0ebc\allowbreak{}5aea26be\allowbreak{}5f5a2193\allowbreak{}db0c32b2} \\
\texttt{tools/\allowbreak{}doc\_\allowbreak{}dirac16\_\allowbreak{}symbols.py} & \texttt{69224cd3\allowbreak{}70206759\allowbreak{}e31abd22\allowbreak{}5fc55663\allowbreak{}0455d6d9\allowbreak{}9226db0b\allowbreak{}c46068ea\allowbreak{}e9b27db0} \\
\texttt{tools/\allowbreak{}doc\_\allowbreak{}gravity.py} & \texttt{50d7c83e\allowbreak{}0232bb23\allowbreak{}b479fc8a\allowbreak{}cca1100a\allowbreak{}8239d781\allowbreak{}8562fece\allowbreak{}3307b1b7\allowbreak{}701ee93a} \\
\texttt{tools/\allowbreak{}doc\_\allowbreak{}gravity\_\allowbreak{}provenance.py} & \texttt{75a4450e\allowbreak{}26a00bf5\allowbreak{}7cf13cab\allowbreak{}5c13a3e1\allowbreak{}177374df\allowbreak{}44cbea3b\allowbreak{}b7a54bf9\allowbreak{}1c6f33dd} \\
\texttt{tools/\allowbreak{}doc\_\allowbreak{}gravity\_\allowbreak{}symbols.py} & \texttt{fcff3b16\allowbreak{}ea23fa62\allowbreak{}fe205e66\allowbreak{}310552f7\allowbreak{}d3e119c6\allowbreak{}19c15037\allowbreak{}9e67bd17\allowbreak{}150b5d22} \\
\texttt{tools/\allowbreak{}doc\_\allowbreak{}ks\_\allowbreak{}gravity.py} & \texttt{9aa73527\allowbreak{}f90318ee\allowbreak{}3444fcf9\allowbreak{}7af41d22\allowbreak{}55a5d8fb\allowbreak{}1697d9fa\allowbreak{}0128ca67\allowbreak{}1e94e51f} \\
\texttt{tools/\allowbreak{}doc\_\allowbreak{}ks\_\allowbreak{}gravity\_\allowbreak{}provenance.py} & \texttt{ae71f9e8\allowbreak{}11bd4277\allowbreak{}32ea9d4f\allowbreak{}b25ddf3e\allowbreak{}7c216df2\allowbreak{}6a2f91ff\allowbreak{}6b9d407a\allowbreak{}264efd16} \\
\texttt{tools/\allowbreak{}doc\_\allowbreak{}ks\_\allowbreak{}gravity\_\allowbreak{}symbols.py} & \texttt{2acdf7f9\allowbreak{}01b9cb7c\allowbreak{}3047b6de\allowbreak{}6895d323\allowbreak{}85be37d1\allowbreak{}6789df9d\allowbreak{}9124f976\allowbreak{}c6d5a515} \\
\texttt{tools/\allowbreak{}doc\_\allowbreak{}provenance.py} & \texttt{ebef465e\allowbreak{}678e11c5\allowbreak{}889047d0\allowbreak{}b4b93191\allowbreak{}c7cf81af\allowbreak{}f9791e75\allowbreak{}c2570b89\allowbreak{}3bd5f7c1} \\
\texttt{tools/\allowbreak{}doc\_\allowbreak{}radion.py} & \texttt{824278aa\allowbreak{}067f4962\allowbreak{}2dcca21d\allowbreak{}5bf19dc1\allowbreak{}ec870b01\allowbreak{}c87d2b2d\allowbreak{}1363e716\allowbreak{}408f2cb1} \\
\texttt{tools/\allowbreak{}doc\_\allowbreak{}radion\_\allowbreak{}common.py} & \texttt{7f4200a4\allowbreak{}7fc921d2\allowbreak{}f537af7e\allowbreak{}9819ecd1\allowbreak{}9a7fbc47\allowbreak{}61335b0c\allowbreak{}980220a9\allowbreak{}9b114068} \\
\texttt{tools/\allowbreak{}doc\_\allowbreak{}radion\_\allowbreak{}provenance.py} & \texttt{5c4d6156\allowbreak{}1df76ea6\allowbreak{}113cf913\allowbreak{}1ad081de\allowbreak{}6ac7aae8\allowbreak{}2a591817\allowbreak{}e7acd670\allowbreak{}c5cf3f52} \\
\texttt{tools/\allowbreak{}doc\_\allowbreak{}radion\_\allowbreak{}symbols.py} & \texttt{fb600386\allowbreak{}675d6000\allowbreak{}ff94373e\allowbreak{}34b442ba\allowbreak{}8ba4f242\allowbreak{}f3692aa3\allowbreak{}99dad84e\allowbreak{}5c90b425} \\
\texttt{tools/\allowbreak{}doc\_\allowbreak{}readme.py} & \texttt{0cb6ed71\allowbreak{}66ec4d08\allowbreak{}f752e827\allowbreak{}87b57915\allowbreak{}27363ea3\allowbreak{}471deccc\allowbreak{}31a05745\allowbreak{}ac29bd52} \\
\texttt{tools/\allowbreak{}doc\_\allowbreak{}readme\_\allowbreak{}stage3.py} & \texttt{af6ffcd0\allowbreak{}8fbde438\allowbreak{}a89d7622\allowbreak{}8b3329b8\allowbreak{}b585c555\allowbreak{}c66cfefe\allowbreak{}6964c487\allowbreak{}055cb0d1} \\
\texttt{tools/\allowbreak{}doc\_\allowbreak{}readme\_\allowbreak{}stage4.py} & \texttt{bd40d0bd\allowbreak{}0fdc2fb9\allowbreak{}b50b453b\allowbreak{}cce7805d\allowbreak{}91c543cd\allowbreak{}594ad92f\allowbreak{}d4beb22e\allowbreak{}76d5765f} \\
\texttt{tools/\allowbreak{}doc\_\allowbreak{}readme\_\allowbreak{}stage5.py} & \texttt{ed2c3a23\allowbreak{}635cca87\allowbreak{}94213032\allowbreak{}81b3a519\allowbreak{}4fb2b21c\allowbreak{}49329b85\allowbreak{}73e6d43e\allowbreak{}cf55f1d4} \\
\texttt{tools/\allowbreak{}doc\_\allowbreak{}readme\_\allowbreak{}stage6.py} & \texttt{8877f882\allowbreak{}2fb22d96\allowbreak{}21686aaa\allowbreak{}b68f215b\allowbreak{}7faedb02\allowbreak{}dc7ad31e\allowbreak{}d20d165e\allowbreak{}162c90af} \\
\texttt{tools/\allowbreak{}doc\_\allowbreak{}readme\_\allowbreak{}stage7.py} & \texttt{59487bf3\allowbreak{}bb233f87\allowbreak{}92d8b626\allowbreak{}728476fc\allowbreak{}fae4d6ff\allowbreak{}1871c1d7\allowbreak{}841d67fa\allowbreak{}c1818bb2} \\
\texttt{tools/\allowbreak{}doc\_\allowbreak{}sm.py} & \texttt{d896153c\allowbreak{}39ba42bf\allowbreak{}67386bd2\allowbreak{}c8deaee4\allowbreak{}fade650b\allowbreak{}d080cf16\allowbreak{}a6bb39d8\allowbreak{}4d29702c} \\
\texttt{tools/\allowbreak{}doc\_\allowbreak{}sm\_\allowbreak{}provenance.py} & \texttt{28f19e04\allowbreak{}777cb914\allowbreak{}9196394f\allowbreak{}c8093fc1\allowbreak{}b67f64d0\allowbreak{}c9cc9756\allowbreak{}0446b3d0\allowbreak{}9e902276} \\
\texttt{tools/\allowbreak{}doc\_\allowbreak{}sm\_\allowbreak{}symbols.py} & \texttt{5ea4c5ec\allowbreak{}d497f34d\allowbreak{}6269a2dd\allowbreak{}3b0aaf80\allowbreak{}90509f52\allowbreak{}73ae89c2\allowbreak{}f111a92a\allowbreak{}ae600a87} \\
\texttt{tools/\allowbreak{}doc\_\allowbreak{}stage3\_\allowbreak{}common.py} & \texttt{981c8418\allowbreak{}03c00c70\allowbreak{}15e49add\allowbreak{}60ec5c55\allowbreak{}f017a294\allowbreak{}4884382e\allowbreak{}812fbd54\allowbreak{}0eede2d5} \\
\texttt{tools/\allowbreak{}doc\_\allowbreak{}status.py} & \texttt{a8f940cb\allowbreak{}0fea5bc8\allowbreak{}8381e77b\allowbreak{}ba1e9c6d\allowbreak{}7ca8992d\allowbreak{}d856688d\allowbreak{}d01620cd\allowbreak{}39cfc2d5} \\
\texttt{tools/\allowbreak{}doc\_\allowbreak{}status\_\allowbreak{}common.py} & \texttt{0e813613\allowbreak{}05921d95\allowbreak{}1b7d614e\allowbreak{}b49eb226\allowbreak{}bd225acc\allowbreak{}a9f30820\allowbreak{}cf80d0d5\allowbreak{}fa71e222} \\
\texttt{tools/\allowbreak{}doc\_\allowbreak{}status\_\allowbreak{}provenance.py} & \texttt{059bbcda\allowbreak{}81f46a9b\allowbreak{}b9a968e0\allowbreak{}72ca6867\allowbreak{}c3e6aaa5\allowbreak{}3940ad14\allowbreak{}75142451\allowbreak{}594eec5b} \\
\texttt{tools/\allowbreak{}doc\_\allowbreak{}status\_\allowbreak{}symbols.py} & \texttt{c9720ccf\allowbreak{}538b5c65\allowbreak{}7bb54ceb\allowbreak{}53e5285b\allowbreak{}005a9fdd\allowbreak{}113dd485\allowbreak{}01862203\allowbreak{}3a6ec1b7} \\
\texttt{tools/\allowbreak{}doc\_\allowbreak{}symbols.py} & \texttt{0d7ba32d\allowbreak{}3213da50\allowbreak{}4dc935d4\allowbreak{}e9930d25\allowbreak{}e37c7294\allowbreak{}ee254d6c\allowbreak{}b3c4aeae\allowbreak{}feee3218} \\
\texttt{tools/\allowbreak{}doc\_\allowbreak{}torus.py} & \texttt{b3045dbe\allowbreak{}9447917b\allowbreak{}d117906d\allowbreak{}3b305f08\allowbreak{}8815d63e\allowbreak{}94439344\allowbreak{}a99265fb\allowbreak{}32d7e6f2} \\
\texttt{tools/\allowbreak{}doc\_\allowbreak{}torus\_\allowbreak{}brane\_\allowbreak{}common.py} & \texttt{43b11a04\allowbreak{}bed8af24\allowbreak{}9206a4fb\allowbreak{}5686678c\allowbreak{}a0a82bf6\allowbreak{}22872eb4\allowbreak{}74644a6e\allowbreak{}b40bfc22} \\
\texttt{tools/\allowbreak{}doc\_\allowbreak{}torus\_\allowbreak{}provenance.py} & \texttt{c9f1e214\allowbreak{}1a14785c\allowbreak{}955a5be1\allowbreak{}d387873f\allowbreak{}52437d39\allowbreak{}c4f23818\allowbreak{}c96d8dbf\allowbreak{}68ffc806} \\
\texttt{tools/\allowbreak{}doc\_\allowbreak{}torus\_\allowbreak{}symbols.py} & \texttt{3327660e\allowbreak{}1f31cc6d\allowbreak{}06c99199\allowbreak{}bc926feb\allowbreak{}4ced37ab\allowbreak{}e27c01a3\allowbreak{}2a3eac57\allowbreak{}0f4a8724} \\
\texttt{tools/\allowbreak{}docgen.py} & \texttt{747d5abe\allowbreak{}53470778\allowbreak{}02cd302a\allowbreak{}08ad293e\allowbreak{}5b458630\allowbreak{}f73ce4f9\allowbreak{}47c06862\allowbreak{}68609ca2} \\
\texttt{tools/\allowbreak{}eval\_\allowbreak{}cells\_\allowbreak{}log.wl} & \texttt{833b6bd8\allowbreak{}54cab253\allowbreak{}e5bb935c\allowbreak{}28f28ab5\allowbreak{}0b286f40\allowbreak{}9d130f43\allowbreak{}3fa5dfb9\allowbreak{}b9236bbe} \\
\texttt{tools/\allowbreak{}extra71.py} & \texttt{9935ef99\allowbreak{}cd9300f2\allowbreak{}a7d4f06d\allowbreak{}36cf7e17\allowbreak{}3912ebae\allowbreak{}61b6feff\allowbreak{}448ebe05\allowbreak{}88853168} \\
\texttt{tools/\allowbreak{}git\_\allowbreak{}push.wl} & \texttt{ac05a6d0\allowbreak{}256834b4\allowbreak{}8cc9fa94\allowbreak{}3fa247bc\allowbreak{}cc574f72\allowbreak{}3e68f053\allowbreak{}d00bfba9\allowbreak{}aca2e21a} \\
\texttt{tools/\allowbreak{}lovelock71.py} & \texttt{d93ecbd0\allowbreak{}f82adb97\allowbreak{}2218d59d\allowbreak{}fe32999a\allowbreak{}4af94494\allowbreak{}5ca97586\allowbreak{}07f3a633\allowbreak{}e1af95a1} \\
\texttt{tools/\allowbreak{}make\_\allowbreak{}7p1\_\allowbreak{}nb.py} & \texttt{996d26c9\allowbreak{}552cd449\allowbreak{}02560af1\allowbreak{}6a1683fa\allowbreak{}71457261\allowbreak{}f18fefae\allowbreak{}ddae87bd\allowbreak{}99a2b144} \\
\texttt{tools/\allowbreak{}make\_\allowbreak{}docs.py} & \texttt{88998b59\allowbreak{}52ce99e0\allowbreak{}e04f65ca\allowbreak{}f9959eef\allowbreak{}de0658e4\allowbreak{}dcfc73d9\allowbreak{}36818e2e\allowbreak{}49c3b400} \\
\texttt{tools/\allowbreak{}make\_\allowbreak{}status\_\allowbreak{}xlsx.py} & \texttt{53e06316\allowbreak{}1fcb9acc\allowbreak{}88749d49\allowbreak{}1c26398e\allowbreak{}499f22da\allowbreak{}46e64792\allowbreak{}a2c828ca\allowbreak{}ecd19a47} \\
\texttt{tools/\allowbreak{}nb\_\allowbreak{}part\_\allowbreak{}coupled.py} & \texttt{e630521d\allowbreak{}64132e71\allowbreak{}afc86111\allowbreak{}4f699eb0\allowbreak{}6454557c\allowbreak{}f14a57ad\allowbreak{}8ec1eef9\allowbreak{}6247aae0} \\
\texttt{tools/\allowbreak{}nb\_\allowbreak{}part\_\allowbreak{}dark.py} & \texttt{10ad8fb8\allowbreak{}ae25fa07\allowbreak{}bbee758c\allowbreak{}653995ac\allowbreak{}d4e2c807\allowbreak{}33059a8b\allowbreak{}3648bd15\allowbreak{}93e46db0} \\
\texttt{tools/\allowbreak{}nb\_\allowbreak{}part\_\allowbreak{}gravity.py} & \texttt{c4caf2ad\allowbreak{}5c932e8f\allowbreak{}603341ed\allowbreak{}ade5f4b9\allowbreak{}8fa784fe\allowbreak{}b05906c0\allowbreak{}51eab7c8\allowbreak{}84cd68f8} \\
\texttt{tools/\allowbreak{}nb\_\allowbreak{}part\_\allowbreak{}ks\_\allowbreak{}gravity.py} & \texttt{a656c0ad\allowbreak{}040b4576\allowbreak{}b75b17e0\allowbreak{}29545109\allowbreak{}e531c8e6\allowbreak{}56ad99bf\allowbreak{}db5cba94\allowbreak{}6e967772} \\
\texttt{tools/\allowbreak{}nb\_\allowbreak{}part\_\allowbreak{}radion.py} & \texttt{53c56047\allowbreak{}07a3e222\allowbreak{}28c840ee\allowbreak{}0a2d2125\allowbreak{}616f9410\allowbreak{}f75a5535\allowbreak{}041dfaf1\allowbreak{}8f497c00} \\
\texttt{tools/\allowbreak{}nb\_\allowbreak{}part\_\allowbreak{}sm.py} & \texttt{43e92bee\allowbreak{}d57453a4\allowbreak{}b57931de\allowbreak{}6d603ac6\allowbreak{}c218efba\allowbreak{}6e400513\allowbreak{}3b69b805\allowbreak{}4a975c5b} \\
\texttt{tools/\allowbreak{}nb\_\allowbreak{}part\_\allowbreak{}torus\_\allowbreak{}brane.py} & \texttt{fb904ff8\allowbreak{}d85b904c\allowbreak{}55b3b2a8\allowbreak{}f60fb0cb\allowbreak{}047f1ccf\allowbreak{}1b98db33\allowbreak{}64092f2b\allowbreak{}ae0b2162} \\
\texttt{tools/\allowbreak{}nbparse.py} & \texttt{ad2255a8\allowbreak{}864be744\allowbreak{}7a550238\allowbreak{}40b327f7\allowbreak{}89cb8c9d\allowbreak{}868d4226\allowbreak{}1b192616\allowbreak{}53847423} \\
\texttt{tools/\allowbreak{}nbwrite.py} & \texttt{dcbe8ce5\allowbreak{}b3203c89\allowbreak{}a1d17fbe\allowbreak{}946e1b6b\allowbreak{}7bb7f889\allowbreak{}5afa9447\allowbreak{}43d651cf\allowbreak{}efadeffb} \\
\texttt{tools/\allowbreak{}octalg.py} & \texttt{36b50120\allowbreak{}b69a963e\allowbreak{}252e07d6\allowbreak{}0ec00775\allowbreak{}5220822b\allowbreak{}4ad33971\allowbreak{}6ea7aed9\allowbreak{}3d58eb24} \\
\texttt{tools/\allowbreak{}ref44.py} & \texttt{ab100c3e\allowbreak{}bcc055fb\allowbreak{}e59c132a\allowbreak{}02f80ecb\allowbreak{}17630e93\allowbreak{}4e99d59f\allowbreak{}88e919ec\allowbreak{}8064faf5} \\
\texttt{tools/\allowbreak{}ref71.py} & \texttt{1a7a06eb\allowbreak{}0add3f45\allowbreak{}9611776d\allowbreak{}c9c474a5\allowbreak{}bb959d94\allowbreak{}7a9db949\allowbreak{}db4612b1\allowbreak{}ab9434e0} \\
\texttt{tools/\allowbreak{}ref\_\allowbreak{}coupled.py} & \texttt{695dbd93\allowbreak{}75753653\allowbreak{}472d9c37\allowbreak{}502f0067\allowbreak{}c4a64fe0\allowbreak{}1a64ba8c\allowbreak{}05e8b2c5\allowbreak{}da85aa58} \\
\texttt{tools/\allowbreak{}ref\_\allowbreak{}dark.py} & \texttt{1600e6b5\allowbreak{}7ab0383b\allowbreak{}4f2b664e\allowbreak{}2796d386\allowbreak{}4c9afa77\allowbreak{}a9ed59f0\allowbreak{}07f66e11\allowbreak{}39453ff5} \\
\texttt{tools/\allowbreak{}ref\_\allowbreak{}dirac16.py} & \texttt{32e8c50d\allowbreak{}e1c480bc\allowbreak{}4ef50d25\allowbreak{}4d04a9aa\allowbreak{}a0e56b48\allowbreak{}6198edf0\allowbreak{}355905c0\allowbreak{}dda25674} \\
\texttt{tools/\allowbreak{}ref\_\allowbreak{}gravity.py} & \texttt{f07cf8ad\allowbreak{}aff585af\allowbreak{}89c87d19\allowbreak{}532f1aaa\allowbreak{}7db606ad\allowbreak{}88ae9048\allowbreak{}3b2d1f41\allowbreak{}341edf36} \\
\texttt{tools/\allowbreak{}ref\_\allowbreak{}ks\_\allowbreak{}gravity.py} & \texttt{b6488ef2\allowbreak{}4b1c16aa\allowbreak{}cd8ff6cf\allowbreak{}6a5a1ba3\allowbreak{}f316e082\allowbreak{}449ad974\allowbreak{}80022d1a\allowbreak{}2638327f} \\
\texttt{tools/\allowbreak{}ref\_\allowbreak{}radion.py} & \texttt{e9d47176\allowbreak{}b9f2d223\allowbreak{}2a84b3b8\allowbreak{}2188bd3b\allowbreak{}0a19cde3\allowbreak{}dbe1ba01\allowbreak{}3c1fe5ab\allowbreak{}ff6341ce} \\
\texttt{tools/\allowbreak{}ref\_\allowbreak{}sm.py} & \texttt{f3f80806\allowbreak{}893f42bb\allowbreak{}f166e1f1\allowbreak{}de896ffd\allowbreak{}611dbfe8\allowbreak{}47685da8\allowbreak{}e3ee367e\allowbreak{}9cde6aba} \\
\texttt{tools/\allowbreak{}ref\_\allowbreak{}torus\_\allowbreak{}brane.py} & \texttt{1f623656\allowbreak{}ab0f5c8e\allowbreak{}e43cb310\allowbreak{}f7bfd580\allowbreak{}cbba5a73\allowbreak{}85da2963\allowbreak{}adaae4d2\allowbreak{}f8bdaaf7} \\
\texttt{tools/\allowbreak{}run\_\allowbreak{}eval.wl} & \texttt{85ba54f1\allowbreak{}8a8362a6\allowbreak{}9c991e66\allowbreak{}a207230d\allowbreak{}d001ba18\allowbreak{}e52639fc\allowbreak{}1edf7cac\allowbreak{}761aac30} \\
\texttt{tools/\allowbreak{}verify\_\allowbreak{}7p1\_\allowbreak{}nb.py} & \texttt{b414dc7b\allowbreak{}69cc03e1\allowbreak{}bc7d4bed\allowbreak{}3a26d836\allowbreak{}7b39e60e\allowbreak{}67ec84ae\allowbreak{}3764a86f\allowbreak{}6bf4c5dd} \\
\texttt{tools/\allowbreak{}watch\_\allowbreak{}requests.wl} & \texttt{d9ceedb0\allowbreak{}b45ac3cf\allowbreak{}9b4dd851\allowbreak{}9ac16495\allowbreak{}3f2c80e5\allowbreak{}188b1557\allowbreak{}fe20ef51\allowbreak{}07ae352d} \\*
\end{longtable}}
\Needspace{12\baselineskip}\section*{5. Process and environment}\addcontentsline{toc}{section}{5. Process and environment}
\begin{itemize}
\item Workspace: \texttt{D:\textbackslash{}\allowbreak{}Developer\textbackslash{}\allowbreak{}github\textbackslash{}\allowbreak{}7+1} on Windows. Python and LaTeX ran in a Linux shell with that folder mounted; intermediate files went to \texttt{build/\allowbreak{}}, which the repository’s .gitignore excludes.
\item The notebook is generated by \texttt{tools/\allowbreak{}make\_\allowbreak{}7p1\_\allowbreak{}nb.py} (sections 0–6 by the code of the first stage, unchanged except for the signature explanation of 27 October 2026; sections 7–10 appended), with cells written by \texttt{tools/\allowbreak{}nbwrite.py} as pure ASCII (named characters such as \texttt{\textbackslash{}\allowbreak{}[CapitalGamma]}, and \texttt{\textbackslash{}\allowbreak{}:03d6} for $\varpi$).
\item All Wolfram Language work ran in the Wolfram 15.0.1 front end (the program of section 4.1), version 15.0.1 for Microsoft Windows (64-bit) (July 2, 2026). Jobs were requested without keyboard input through \texttt{tools/\allowbreak{}watch\_\allowbreak{}requests.wl}: loaded once in that front end, it reads \texttt{build/\allowbreak{}request.txt} every 10 seconds and evaluates \texttt{tools/\allowbreak{}run\_\allowbreak{}eval.wl} (job \texttt{eval}), \texttt{tools/\allowbreak{}git\_\allowbreak{}push.wl} (job \texttt{push}) or a development file such as \texttt{tools/\allowbreak{}eval\_\allowbreak{}cells\_\allowbreak{}log.wl} (job \texttt{run}).
\item Python 3.13.17 with sympy 1.14.0, mpmath 1.3.0 and numpy 2.5.3; fontTools reads the font coverage for the LaTeX output.
\item This record and the main document and symbol dictionary of this stage are generated by \texttt{tools/\allowbreak{}make\_\allowbreak{}docs.py} from the modules \texttt{docgen.py}, \texttt{doc\_\allowbreak{}common.py}, \texttt{doc\_\allowbreak{}dirac16\_\allowbreak{}common.py}, \texttt{doc\_\allowbreak{}dirac16.py}, \texttt{doc\_\allowbreak{}dirac16\_\allowbreak{}symbols.py} and \texttt{doc\_\allowbreak{}dirac16\_\allowbreak{}provenance.py} of the folder tools/; one source renders Markdown and LaTeX, and every matrix, check list and number is computed by Python or parsed from the evaluated 7+1.nb. The LaTeX is compiled with XeLaTeX (TeX Live 2022, \texttt{latexmk -xelatex}, fonts Latin Modern and DejaVu).
\item Pushes: \texttt{tools/\allowbreak{}git\_\allowbreak{}push.wl} runs the user’s Git for Windows from the Wolfram front end with \texttt{GIT\_\allowbreak{}TERMINAL\_\allowbreak{}PROMPT=0} and \texttt{credential.interactive=never} (no login prompt), pushes the branch main to origin \texttt{https:/\allowbreak{}/\allowbreak{}github.com/\allowbreak{}once-ere/\allowbreak{}7plus1.git}, and logs \texttt{git ls-remote origin}; the published commit is compared with the local one after every push.
\item Order of the work: backup of 7+1.nb; Python reference \texttt{tools/\allowbreak{}ref\_\allowbreak{}dirac16.py}; notebook sections 7–10, developed cell by cell; evaluation and cross-check (\texttt{tools/\allowbreak{}verify\_\allowbreak{}7p1\_\allowbreak{}nb.py}); documentation. The notation of the complex scalar source was changed from $\mu$ to $\varpi$ after the first evaluation, and 7+1.nb was regenerated and evaluated again; the boost, the proof matrices and Euler’s Gamma function were renamed in the documents only ($\mathcal B$, $Y$, $Z$, $\Gamma_{\rm E}$).
\end{itemize}
\Needspace{12\baselineskip}\section*{6. Authorship}\addcontentsline{toc}{section}{6. Authorship}
Notebook sections 7–10, tools and documents written by Claude (Anthropic, model claude-opus-5-5) for Nash, completed on 2026-10-26 (this record was regenerated on 2026-10-28), from the task described in section 1 and the first part of 7+1.nb described in section 2.1.

\end{document}
